\subsection{Absolute convergence, support at one split place, and the inert quotient}
\label{threal:main}

Write \(d=a+a'\), \(f=[F:\mathbb Q]\), and use the height \(H\) of
\S\ref{newtheta:global}. Put \([G]=G(F)\backslash G(\mathbb A_F)\).
There is a difference between the estimates needed to define the two
periods and the identity comparing their values. We prove the estimates
for the undoubled groups.

\begin{proposition}\label{threal:decay}
For the signatures and ranks of \S\ref{newtheta:global}, the unprimed
theta kernel is absolutely integrable over \([U(V)]\) at every \(h\)
for every adelic Schwartz function. For a primed test \(\Phi'\) in the
specified full-rank space at \(v_1\) and every \(N\geq0\), there exist
\(C_{\Phi',N}\) and an integer \(k\geq0\) such that
\[
 |\Theta_{\Phi'}(g,h)|\,H_{[U(V')]}(g)^N
                      \leq C_{\Phi',N} H(h)^k .
\]
The primed kernel is absolutely integrable in \(g\), its integral has
moderate growth in \(h\), and the full-rank test space is invariant
under the two unitary actions.
\end{proposition}
\begin{proof}
Use \(Y\) for either source, \(m=\dim_EY\), and let \(r\) be its
\(F\)-rational Witt index. Fix rational hyperbolic pairs
\((e_i,e_i^*)_{1\le i\le r}\) and their orthogonal complement \(Y_0\).
In \(Y\otimes_EW\), take one whole half of each hyperbolic pair
tensored with \(W\) as position coordinates and complete by a rational
Lagrangian in the underlying symplectic space of \(Y_0\otimes_E W\).
The theta sum now has positions
\[
        (x_1,\ldots,x_r,y)\in W(E)^r\times F^b,\qquad
             b=(m-2r)d ,
\]
where \(W(E)\) denotes the underlying \(E\)-vector space.
Choose the half for which the torus in the positive chamber acts on
functions by
\[
 \phi(x,y)\longmapsto
       \left(\prod_i T_i^{d/2}\right)
       \phi(t_1x_1,\ldots,t_rx_r,y),\qquad
          T_i=t_i^{2f},\quad t_i>0 ,
 \tag{*}
\]
up to a scalar of absolute value one. Indeed each of the \(d\)
\(E\)-coordinates contributes \(2f\) real dimensions and the
prefactor is the square root of the real change of volume. The torus
has value \(t_i\) at all real places and unit finite components.
With the parabolic stabilizing the expanding isotropic flag one uses
the opposite position half: the action on arguments is inverse to
the action on translations.

The change from the given Darboux coordinates is rational. Indeed the
images of a rational basis are rational by assumption, and their
determinant is nonzero. Rational linear changes, rational quadratic
phases and partial Fourier transforms preserve the absolute value of
the rational summation functional, by the product formula, triviality
of \(\psi\) on \(F\), and Poisson summation with self-dual measures.
These operations generate the rational symplectic group by the Siegel
Bruhat decomposition. They also show that the absolute value of the
kernel is invariant under rational left translates in either group.
These arguments apply to each allowed unitary implementation: two
unitary operators implementing the same action on the Heisenberg
operators differ by a scalar of modulus one. In Schr\"odinger
coordinates their quotient commutes with every modulation and hence
is multiplication on \(L^2\), and commuting with the translations
makes that multiplier constant. This argument also permits the
use of the usual tensor-product implementations for estimates at a
fixed adelic group element.

Apply adelic reduction theory to the connected reductive group \(U(Y)\).
Its center is the norm-one torus and has no \(F\)-split part.
A minimal \(F\)-parabolic is the stabilizer of the chosen flag.
The quotient is covered by a Siegel set with representatives
\(n m_0 t k_0\), with \(n,m_0,k_0\) in fixed compact sets of the
unipotent radical, norm-one part of the centralizing Levi, and
maximal compact subgroup, respectively. Here \(t_i/t_{i+1}\) and
\(t_r\) are bounded below by a positive constant. When \(r=0\) the
set is compact. These are the compactness and fundamental-set theorems
in adelic reduction theory, with inversion to pass from the right
quotient to the left quotient \cite[\S\S2.1--2.2]{springerReduction}.
The elements \(t^{-1}nt\) stay in a compact set, since the inverse
of every positive root is bounded in this chamber. Thus the
representatives have the form \(tc\), with \(c\) in a fixed compact
set. In the integration coordinates above the Haar density in \(t\)
is a constant times
\[
       \prod_{i=1}^{r} T_i^{-(m-2i+1)}\,\frac{dT_i}{T_i}.
\]
Indeed for each \(j>i\) the two relative roots with weights
\(t_i/t_j,t_it_j\) have parameter in \(E\); the root \(t_i\) has
parameters in the \(m-2r\) dimensional space \(Y_0\) over \(E\);
the root \(t_i^2\) has one parameter over \(F\). Their determinant
on the unipotent radical is
\(\prod_i T_i^{2(r-i)+(m-2r)+1}\).
Integration of a nonnegative rationally invariant function over a
covering Siegel set bounds its quotient integral from above.

Here are the bounds on the test in (*) as \(c,h\) vary. Fix coordinates
over \(\mathbb Q\) using bases of the number fields and put \(D=fmd\).
For a pure archimedean and finite tensor, the finite part of
\(\omega(c,h)\Phi\) is supported on \(M^{-1}\widehat{\mathbb Z}^{D}\)
and has supremum at most \(C M^{D/2}\), for an integer
\[
             1\le M\le C_\Phi H(h)^A .
\]
To see this, start with an integer \(M_0\) such that translations and
modulations by \(M_0\widehat{\mathbb Z}^{2D}\) fix the finite test,
using the standard \(\mathbb Q\)-adelic character and its dual
coordinates. Increasing \(M_0\) also fixes coordinate denominators.
After conjugation it suffices to have
\((c,h)^{-1}M\widehat{\mathbb Z}^{2D}\subset
 M_0\widehat{\mathbb Z}^{2D}\) in the symplectic representation.
At each place \(v\), put \(n_v\) equal to the largest denominator
order of the entries of \(h,h^{-1}\), or zero if all are integral,
using all factors of \(E\otimes F_v\). The extension \(E/F\) is
unramified at finite places. In fixed bases the entries on tensors
involve only matrix multiplication and fixed changes of coordinates.
All the \(c\)'s carry bounded denominators, supported at finitely many
primes. Thus \(M=C_0\prod_{p}p^{\max_{v\mid p}n_v}\) suffices for
a fixed positive integer \(C_0\). Since \(p\le q_v\) and the finite
height factor is at least \(q_v^{n_v}\), this proves the bound.
Invariance by modulations forces the stated support; invariance by
translations makes the function constant on cosets of
\(M\widehat{\mathbb Z}^{D}\), each of volume \(M^{-D}\) in product
measure of integral volume one. Unitarity gives the supremum bound.

At infinity, for every \(B\ge0\), the absolute value of the transformed
test at \(z\) is at most \(C_B H(h)^{A_B}(1+\|z\|)^{-B}\).
Indeed commute multiplication and differentiation through the unitary
operators using Heisenberg covariance. Each is a linear combination of
the original position and momentum operators, with coefficients
bounded by a fixed power of the maximum of the archimedean matrix
norms of \(h,h^{-1}\) and a constant for \(c\). Apply this identity to
words of length at most \(B+D+2\) (round \(B\) up to an integer).
Euclidean Sobolev embedding applied to each monomial times the
transformed function then gives the assertion. Finite sums of pure
tensors satisfy the same majorant. In these estimates \(A_B\) may
depend on the fixed spaces, \(B\), and the test. The specified height
uses the local spectral absolute values over \(F_v\), extending
\(|\pi_v|=q_v^{-1}\), and takes the maximum of \(1\) and the norms of
both matrices.
\src{src/Development/AbelianVarietyTheory/Definitions/Foundations/EventuallyNormFinLeOne.lean}{405}{Matrix height across the places}.

For the unprimed integral fix \(h\); then \(M\) is bounded and
\(t_i\gg1\). Summing the Schwartz majorant with \(B>D+1\) on
\(M^{-1}\mathbb Z^D\) gives \(O_{\Phi,h}(\prod_i T_i^{d/2})\).
For \(r\le p\) every exponent \(d/2-(a-2i+1)\) is strictly negative,
since \(a-d/2=(a-a')/2>2(p+q)+1\). Integration even over the larger
rectangle \(T_i\ge c_1>0\) converges.

For the primed integral we verify the support after the change of
polarization. At \(v_1\) take the specified conjugate split coordinates,
so the original position variable is an \(m\)-by-\(d\) matrix over
\(F_{v_1}\). The two unitary actions there are invertible row and
column changes times a scalar, since they preserve each of the two
split Lagrangians. Therefore they preserve full row rank. This also
proves invariance of the adelic span by factorization, since all other
factors stay Schwartz and a restricted adelic element fixes the
standard lattice factors up to scalar almost everywhere. For clarity,
write the conjugate split matrix first in the rational Witt row basis.
In a hyperbolic pair the two position rows are indexed by \(e_i,e_i^*\).
In the new model the \(E\otimes F_{v_1}\)-valued position block \(x_i\)
contains \emph{one unchanged row} and the Fourier dual of its paired
row. Indeed the tensor symplectic pairing pairs the \(e_i\)-row on
one split factor with the \(e_i^*\)-row on the other and conversely.
One exchanges the position and momentum in one of these two row
pairs, using an invertible linear identification with \(W\) from
its form. The remaining change of polarization takes place in the
\(Y_0\)-block. If the unchanged row vanishes, the original matrix has
deficient rank for every value integrated over in these partial
Fourier transforms. The transformed function therefore vanishes when
\(x_i=0\). This holds separately for each \(i\) and survives finite
sums. The coordinates called for in the support definition are unique
up to scalar on Schwartz--Bruhat functions: an intertwiner with the
stated translation and modulation laws is unique by the local
Schr\"odinger commutant. Algebraically this follows by averaging
modulations over compact open subgroups to project onto support in
any ball, and then translating the indicators of balls.

For a surviving rational summand \(\|x_i\|\ge c_2/M\): the rational
points in the indicated finite support lie in \(M^{-1}\mathbb Z^D\),
and the coordinates were chosen block by block. Given \(L\ge0\),
spend \(rL\) powers in the archimedean majorant to obtain
\(\prod_i (1+t_i\|x_i\|)^{-L}\ll M^{rL}\prod_i t_i^{-L}\).
Another \(D+2\) powers suffice to sum the remaining majorant over
\(M^{-1}\mathbb Z^D\), with bound \(O(M^D)\) by unit-cube counting.
Thus, for every \(J\ge0\),
\[
             |\Theta_{\Phi'}(tc,h)|
                    \ll_{\Phi',J} H(h)^{A_J}\prod_i T_i^{-J}.
\]
Here \(L\) was chosen to absorb both \(2fJ\) and the prefactor in (*).
The quotient height at any translate is at most \(H(tc)\), which on
this set is bounded by \(C\prod_i T_i^{A_0}\), by the matrix formula
for \(H\). Taking \(J\ge NA_0\) gives the claimed inequality.
Integrating the bound for any sufficiently large \(J\) proves
absolute convergence and moderate growth, including the case \(r=0\)
with empty products. All bounds use absolute values, so they apply
to the prescribed Haar fundamental domains.
\end{proof}

\begin{corollary}\label{threal:limits}
An identity with one fixed scalar for all polynomially matched pairs
whose primed input satisfies the support condition implies the same
identity for every matched pair in \S\ref{newtheta:global}.
\end{corollary}
\begin{proof}
Fix \(h\) and a finite Schwartz factor. The unprimed proof bounded the
absolute integral by finitely many archimedean Schwartz seminorms.
The same is true on the primed side for a fixed admissible finite
factor: the integer \(M\), the size of the finite factor and the
compact set used in reduction are fixed, and the Sobolev argument
uses finitely many seminorms of the varying archimedean factor.
For the primed approximations the individual finite factors in the
definition of matching need not themselves satisfy the support
condition. Let \(B\) be their finite-dimensional span and \(B^\circ\)
the subspace of the full finite Schwartz space spanned by admissible
finite tensors. For \(S=\mathcal S(\mathbb R^D)\) one has
\[
      (S\otimes B)\cap(S\otimes B^\circ)=S\otimes(B\cap B^\circ)
\]
inside the algebraic tensor product with all finite Schwartz tests,
as follows by extending a basis of \(B\cap B^\circ\) in each factor
space \(B,B^\circ\). Products of functions realize this tensor product
faithfully: evaluations span the dual of every finite-dimensional
space of functions. Express the approximants and their limit in a
basis of \(B\) adapted to \(B\cap B^\circ\). The coefficients outside
the intersection vanish and the others converge in \(S\), being fixed
finite linear combinations of the original coefficients. Continuity
on each resulting finite slice permits passage to the limit at \(h\).
\end{proof}

\begin{proposition}[Even source ranks at an inert place]\label{threal:inert}
Suppose \(v\) is inert and \(a=2u,\ a'=2w\). The two local
Schwartz--Bruhat spaces have a common irreducible \(U(W)(F_v)\)
quotient, invariant under their respective source groups. The maps
can be chosen to send the prescribed vacua to the same nonzero
spherical vector. Its unramified parameter, in the coordinates
\(q^{2\sigma_i}\) up to permutation and separate inversions, is
\[
      (q,q^3,\ldots,q^{2u-1})\ \sqcup\
      (q,q^3,\ldots,q^{2w-1}),\qquad q=|\mathcal O_F/v|.
\]
\end{proposition}
\begin{proof}
Put \(k=F_v,\ l=E_{\mathfrak w}\), and \(R=u+w\). This is an
unramified quadratic extension of odd residual characteristic, \(q>3\),
and \(\psi_v\) has conductor \(\mathcal O_k\). Since all ramification
of \(\alpha,\beta\) is split, their local components are unramified.
Both are trivial on \(k^\times\) by evenness of the ranks. A
uniformizer of \(k\) is also a uniformizer of \(l\), so both local
characters, and both splitting pairs, are trivial.

The two source lattices are perfect Hermitian lattices. The absolute
trace-dual condition on the target also gives a perfect skew-Hermitian
lattice here. Indeed the standard global additive character with the
specified infinite component has local annihilator of \(\mathcal O_k\)
equal to the absolute codifferent; this follows from the trace
pairing and the uniqueness of the character on \(\mathbb A_F/F\)
with that infinite component (the dual is \(F\)). The conductor
hypothesis makes this codifferent the unit ideal, and \(l/k\) adds
no different. A trace-zero unit converts a perfect skew-Hermitian
form into a Hermitian one. Such a unit is obtained by subtracting
the conjugate of a lift of a residue element outside the base field.
Every even-rank perfect Hermitian lattice here is integrally
hyperbolic. Over the residue fields diagonalize and use surjectivity
of the norm to obtain a nonzero isotropic vector in dimension at
least two. The derivative of the norm form there is the surjective
residue trace pairing with this vector. Hensel lifting gives a
primitive isotropic \(x\). Choose \(y\) pairing to one with \(x\)
and replace \(y\) by \(y-h(y,y)x/2\). The resulting hyperbolic plane
splits integrally, giving induction. In particular the target admits
the integral hyperbolic polarization and both source groups have the
quasi-split sign \(+\).

Use the ordinary Schr\"odinger model on \(Y^R\), for \(Y\) one of the
two Hermitian sources. The lattice vector is, up to scalar,
\(\mathbf1_{\Lambda_Y^R}\). To verify the identification with the
prescribed vacuum, apply its Heisenberg eigenline condition on each
of the two tensor lattice halves of the hyperbolic polarization of
\(W\). The quadratic phase is one on each half, which is isotropic
over \(l\) and stable under the trace-one element. The two conditions
are invariance under the translations by \(\Lambda_Y^R\) and the
modulations by its dual. They force support in \(\Lambda_Y^R\) and
constancy there. The change to these coordinates is a unitary
Heisenberg intertwiner and the original vacuum is nonzero.

We recall why the splitting, including its scalar character, agrees
with the one in the spherical calculation. The prescribed splitting
on each factor extends to its hyperbolic doubling. In the doubled
diagonal polarization evaluation at zero has on its Siegel stabilizer
the law \(\chi(\det t)|\det t|_l^{\dim_EY/2}\), with the oscillator
partner's dimension in the exponent if it is the source that is
doubled. This is precisely the standard doubled splitting with
character \(\chi\). To use the adelic law at a single place, take
evaluation in the diagonal polarization at all places. This is an
admissible functional: rational changes to that polarization are
compositions of the linear changes, phases and Fourier transforms
above, and evaluation is continuous on each fixed finite Schwartz
slice. The functional is nonzero. After unitary identification of the
doubled Heisenberg implementations, their quotient on elements
supported at \(v\) is a scalar character by the commutant argument.
The common evaluation law makes it trivial on the Siegel stabilizer.
Its conjugates generate the hyperbolic unitary group: the Siegel
Bruhat double cosets are represented by interchanging some of the
hyperbolic pairs. In each hyperbolic plane the interchange, up to
a diagonal element, is the product of three alternating upper and
lower root unipotents by the \(2\)-by-\(2\) identity in
\(\mathrm{SL}_2(k)\simeq SU(1,1)(k)\). Upper unipotents are in the
stabilizer and lower ones in its opposite conjugate. Thus the
character is trivial. Restrict on block elements \((g,1)\) of the
doubling and apply the defining product rule on pure tensors to
recover the given single-factor action. This argument works on both
factors. In particular it excludes a scalar character twist that
Heisenberg covariance alone would permit.
\src{src/Development/AbelianVarietyTheory/Definitions/Divisors/HilbertNinetyQuotientHomEqIff.lean}{435}{Doubled Siegel and block restriction laws}.

In Liu's notation take the larger skew-Hermitian group \(G_R\) and
the Hermitian space of dimension \(2s\), first \(s=u\), then \(s=w\).
His unramified additive character and trivial splitting character are
the ones just checked (a unit scaling in the trace pairing changes
only the choice of conductor-zero additive character or hyperbolic
basis). The trivial representation of \(U(2s)\) has parameters
\((q,q^3,\ldots,q^{2s-1})\), since the exponents of the normalized
Borel modulus on its \(l^\times\)-coordinates are
\(s-i+1/2\), for \(1\le i\le s\). In fact the root computation in
the preceding proof gives the modulus before taking its square root.
Liu's Theorem~6.2(2), Proposition~3.4 and the parameter maps of
\S2 say that, for \(R\ge s\) and sign \(+\), the small lift of the
spherical representation on \(G_R\) with the \(s\) parameters inverted
and \(q,q^3,\ldots,q^{2(R-s)-1}\) appended is the spherical
representation with those \(s\) parameters on \(U(2s)\).
The small-lift quotient sends \(\mathbf1_{\Lambda_Y^R}\) to a
nonzero vector, also explicitly in the proof of his Theorem~1.1.
The two applications give the displayed identical Weyl orbit.
Accordingly the oscillator has in each case a quotient onto
\(\pi\boxtimes{\bf1}\) with the same irreducible spherical \(\pi\)
on \(G_R\). The image of each vacuum belongs to the one-dimensional
hyperspecial fixed line; rescale one map to give equality. If maps
on the full space of all functions are requested, extend linearly
from Schwartz--Bruhat functions using any vector-space complement:
the conditions concern only Schwartz--Bruhat vectors.
For these choices of spherical vector, parameter and normalized
induction see \cite[\S\S2, 3 and 6]{liuThetaInert}.
\end{proof}

\begin{proposition}[All inert places, including the odd row ranks]
\label{threal:allinert}
At every inert place, in either parity, the two local Schwartz spaces
have a common irreducible target quotient annihilating the source
differences. The prescribed vacua have the same nonzero image after
one scalar normalization.
\end{proposition}
\begin{proof}
Write \(d=2R\), \(Y=V\) or \(V'\), \(m=\dim_EY\), and
\(\xi=\beta_v\) or \(\alpha_v\), respectively. These last two
characters of \(E_v^\times\) are the same: both are unramified
and their values at a uniformizer of \(F_v\), also a uniformizer
of \(E_v\), are \((-1)^m\). Thus \(\xi^2=1\). On the source the
splitting character is trivial in both cases. Use the integral
hyperbolic polarization of \(W\) in the preceding proof. In this
model \(U(Y)\) acts by linear substitution on \(Y^R\): in the
Siegel Levi implementing its action the determinant modulus is
one, with trivial splitting character. Evaluation at zero is
source invariant. Consequently
\[
        \phi\longmapsto
          \bigl(h\longmapsto(\omega(1,h)\phi)(0)\bigr)
\]
gives a nonzero equivariant map for the target into the normalized
Siegel induction \(I(s_m,\xi)\) of \(U(2R)(F_v)\) from
\((\xi\circ\det)|\det|_{E_v}^{s_m}\),
where \(s_m=(m-R)/2\). Indeed the Siegel unipotent acts by a
quadratic character taking value one at zero. The Levi acts at
zero by \(\xi(\det b)|\det b|_{E_v}^{m/2}\), and the half modulus
of this Siegel parabolic is \(|\det b|_{E_v}^{R/2}\).
The image is nonzero already on the lattice indicator at \(h=1\).
These substitutions are the standard Siegel splitting of the
polarized Weil representation. The doubled uniqueness argument
above identifies them with the actions specified here.

The image has finite length as a submodule of parabolic induction
of an admissible finite-length representation. Choose any
irreducible quotient of this nonzero image. The composite is a
nonzero map from the oscillator onto the trivial source times
this quotient. The unramified vacuum theorem says that such a map
takes the lattice vacuum to a nonzero vector in the hyperspecial
fixed line when the source representation is spherical. Its
hypotheses are precisely odd residue characteristic, an unramified
quadratic extension and additive character of conductor \(\mathcal
O_v\), self-dual lattices and unramified splitting characters. To
interchange the Hermitian and skew-Hermitian spaces one can use a
trace-zero unit and its inverse. The source vacuum in the given
coordinates is a nonzero multiple of this lattice vector by the
Heisenberg characterization above.
\cite[Appendix, Theorem~2.A.1]{harrisSpherical}

In particular this irreducible quotient has the spherical Hecke
character of \(I(s_m,\xi)\): hyperspecial invariants are exact by
compact averaging, and the induction has a single fixed line by
Iwasawa decomposition. For \(m'=2R-m\) one has \(s_{m'}=-s_m\).
The two inducing data are Weyl conjugate on their unramified
parameters: the long Siegel Weyl element takes \(b\) to its
inverse conjugate transpose, \(\xi\) is conjugation invariant
and self-inverse, and the \(R\) exponents of the trivial
\(\mathrm{GL}_R\) representation are symmetric about zero.
The Satake classification thus identifies the two spherical
irreducibles. Scaling the maps to their common fixed line gives
the desired normalization.
\end{proof}

\begin{proposition}[Ordinary relatively split places]
\label{threal:ordinarysplit}
At every place split in \(E/F\) outside \(S=S_0\sqcup Q\), there
is a common irreducible quotient with matching nonzero vacua,
including when the underlying rational prime does not split
completely in \(E\).
\end{proposition}
\begin{proof}
Use the two row and dual-row models of
Proposition~\ref{newtheta:split}. The relative splitting
\(E\otimes F_v=k\times k\) suffices for these models. In either
polarization the two idempotent summands of the tensor lattice give
translations by a full lattice and modulations by its character dual.
Indeed the source lattice is perfect for its \(E\)-valued form,
and the partner lattice is perfect for the absolute trace pairing.
Thus, in dual bases, the modulation coefficients range over the
inverse different of \(k/\mathbb Q_\ell\) times the ordinary dual
of the position lattice. This gives exactly the character dual for
the trace additive character. The latter is forced by its infinite
component and triviality on \(F\), since \(F\) is dense in its
finite adeles. On each idempotent summand separately the quadratic
phase in the vacuum condition is one by isotropy, also in residue
characteristic two. Translation invariance and modulation invariance
force any nonzero vacuum to be a constant times the indicator of
the position lattice. Choose position lattice bases and the
conjugate dual bases so its indicator is the integral-matrix test
on each side and the two target compact subgroups agree. Scaling
these local coordinate intertwiners identifies the given vacua
with the two indicators. The unitary linear action is the prescribed
one by the block and Siegel laws used above; conjugating to the dual
polarization sends the target action to inverse transpose as in
\eqref{newtheta:big}.

In the notation of Proposition~\ref{split:good} put \(n=a\),
\(t=\nu(\pi)\), \(s=d/2\) and
\(D_n(z)=\prod_{j=0}^{n-1}(1-tq^{j-z})\).
For any matrix test \(H\), the cleared determinant integral
\(D_n(z)Z_n(H,z)\) is a Laurent polynomial in \(q^{-z}\).
Here is the refinement needed when \(t\) has finite order. Keep the
integrand exponent \(z-n\) fixed as in that proposition. In size
\(j\) a ball about zero scales to \(M_j(\mathcal O_v)\). Its
denominator divides
\(\prod_{i=0}^{j-1}(1-tq^{n-j+i-z})\) by the same Hermite
count with parameter \(z-n+j\) and the conversion from additive
to group Haar measure. Its numerator is a constant times a
Laurent monomial. Subtract the value at zero times such a ball.
The remaining compact support is covered by pivot charts; the
Schur-complement argument reduces to size \(j-1\), whose denominator
is a subproduct of the one just displayed. Finite disjoint
subdivision and induction prove the claim starting with size zero.

Multiply the larger transform by \(D_n(z)\) and evaluate at \(s\).
The source and parabolic covariance pass from the right half-plane
to this value by the meromorphic identity principle, including if
the value is a residue. The character and modulus factors in that
covariance are holomorphic in \(z\). Smoothness in the target
follows from the compact open stabilizer of the test. The larger
vacuum has value one at the identity by the exact Hermite count,
with \(\operatorname{vol}(\mathrm{GL}_n(\mathcal O_v))=1\).
The smaller vacuum also has nonzero value there, since \(a'<s\),
and can be normalized to one. The two spherical vectors are now
equal. The common induced target of Proposition~\ref{newtheta:split}
is irreducible by Bernstein's theorem, so both nonzero maps are
onto. The stronger assertion on the entire source-difference
kernel at totally split good primes was proved separately in
Proposition~\ref{split:good}.
\end{proof}

\begin{proposition}[The real quotient for \(p=q=0\)]\label{threal:compact}
At each infinite place write \(A=m_\alpha,\ B=m_\beta\) and
\[
        j=(A-B-d)/2,\qquad l=(B+a)/2 .
\]
The two real polynomial modules have the same source-coinvariant
representation of the compact target \(U(d)\), with dominant highest
weight \(((l+j)^a,l^{a'})\). Superscripts indicate repetitions.
The tensor product of these representations over the real places
realizes the real matching maps for this branch.
\end{proposition}
\begin{proof}
The restrictions to \(\mathbb R^\times\) give \(A\equiv a'\),
\(B\equiv a\pmod2\). Thus \(j,l\) are integers, and \(j>0\) by the
weight condition. Let \(D_0=\mathbb C^d\), \(U_0=\mathbb C^a\),
\(U_1=\mathbb C^{a'}\) with the standard actions. In orthonormal
column frames the operators on the polynomial Fock modules are the
following two representations of source times target:
\[
\begin{split}
 &(\det_{U_0}^{-j}\boxtimes\det_{D_0}^{\,l})
               \otimes\operatorname{Sym}(U_0\otimes D_0),\\
 &(\det_{U_1}^{\,j}\boxtimes\det_{D_0}^{\,l+j})
               \otimes\operatorname{Sym}(U_1^*\otimes D_0^*).
\end{split}
\]
We verify the weights to fix the sign convention. In one unprimed
complex matrix entry the tensor form is
\(-\operatorname{Tr}_{\mathbb C/\mathbb R}(-i z\bar z')
=2\,\operatorname{Im}(\bar z z')\); it changes sign on the primed
side. The additive character is \(x\mapsto e^{-2\pi i x}\). In the
positive plane the Heisenberg derivatives for \(1,i\) can be taken
as \(\partial_x,-4\pi i x\). Acting on
\(\exp(-2\pi x^2)\), the second derivative operator is \(i\) times
the first. Thus a rotation acts on the creation operator
\(d\rho(1)-i\,d\rho(i)\) by its ordinary character. Reversing the
form gives the conjugate character. The infinitesimal metaplectic
rotation on the positive plane has eigenvalues \(i(n+1/2)\):
in the creation and annihilation operators \(b^\dagger,b\) with
\([b,b^\dagger]=1\) its generator is the symmetrized quadratic
\(i(b^\dagger b+bb^\dagger)/2\); the negative plane reverses this
sign. Therefore the metaplectic vacuum has half the sum of the signed
entry weights.

The doubled Siegel normalization contributes the additional half of
the splitting-character exponent on the determinant of each single
unitary factor. Indeed on the doubled Siegel Levi \(t\in
\mathrm{GL}_n(\mathbb C)\) the action of the metaplectic Lie algebra
on evaluation gives the positive modulus; its infinitesimal phase
is zero, as one also sees from the real change of coordinates.
The prescribed additional phase on this Levi is
\((\det t/|\det t|)^m\). The determinant of the doubled unitary
element on the same Levi is \(\det t/\overline{\det t}\); hence the
difference of the two infinitesimal implementations is
\((m/2)\operatorname{tr}\). They differ by a Lie algebra character,
and every such character on \(\mathfrak u(n,n)\) is a multiple of
trace since its derived algebra is \(\mathfrak{su}(n,n)\).
Restriction to the first block preserves the trace. Consequently the
source vacuum powers are \((d+B-A)/2=-j\) and
\((-d+A-B)/2=j\); the target powers are \((a+B)/2=l\) and
\((-a'+A)/2=l+j\). The unitary groups are connected, so the
infinitesimal action determines these integrated representations.
The Gaussian times the Hermite polynomials form the polynomial
module here: the annihilation equations for the decaying Gaussian
separate into the one-coordinate first-order equations above, with
the opposite choice of sign giving an increasing exponential, and
the Heisenberg derivatives generate all the creation monomials.

Use the Cauchy decomposition
\(\operatorname{Sym}(U\otimes D_0)=
\bigoplus_\lambda S_\lambda U\otimes S_\lambda D_0\);
this follows by applying Schur--Weyl duality to each homogeneous
degree and taking simultaneous permutation invariants.
The only partition for which \(S_\lambda U_0=\det_{U_0}^j\) is the
rectangle \((j^a)\), with multiplicity one. On the dual side the
rectangle is \((j^{a'})\). Thus the first invariant space is
\(\det^l S_{(j^a)}D_0\), and the second is
\(\det^{l+j} S_{(j^{a'})} D_0^*\). Their highest weights agree,
since duality reverses and negates the list of weights. Both spaces
are irreducible. Averaging over the source compact group projects
onto the invariants and kills precisely the span of the differences
on the polynomial module, by complete reducibility on each finite
degree. Composing with an identification of the two highest-weight
modules gives the claimed surjections. Their tensor product is
irreducible for the product of the compact targets. The maps may
be extended linearly off the polynomial spaces, since the
specified infinitesimal and compact transformation laws concern
polynomial inputs and these actions preserve that subspace.
For comparison with the source conventions see
\src{src/Development/WeilBound/ThetaLifts/LocalTheory/ToralCohomology/PiInfUnitNe.lean}{116}{The unprimed target highest weight}.
\end{proof}

\begin{proposition}[The real quotients for all signatures]\label{threal:allreal}
At each infinite place let \(A=m_\alpha,\ B=m_\beta\).
In Paul--Trapa's normalization the common real representation
is the cohomologically induced module on the ordered Levi
\[
 U(a-p,p)\times U(a'-q,q),\qquad
 A_{\mathfrak q}(\det^{A/2}\boxtimes\det^{B/2}).
\]
The half-determinant notation uses the covers in their normalization.
The output after twisting is an honest module for the target Lie
algebra and its maximal compact group.
\end{proposition}
\begin{proof}
Put \(r=d-p-q,\ t=p+q,\ k=A-B>d\). Work first with the
metaplectic oscillator, before the two splitting twists. With the
additive character and tensor form in Proposition~\ref{threal:compact},
a positive source row has vacuum determinant exponent \((r-t)/2\);
a negative source row has the opposite exponent. This follows by
summing the signed one-entry weights \(1/2\) there. It is the
convention of Paul--Trapa's joint harmonics, their formula (3.3).
In passing to the prescribed actions the source determinant exponents
are shifted by \((B-A)/2\) and \((A-B)/2\), respectively.
The doubled trace argument for that shift in
Proposition~\ref{threal:compact} applies to each signature, since
the doubled group and its Siegel Levi have the same trace character.
Thus the prescribed trivial source corresponds in the metaplectic
oscillators to \(\det^{k/2}\) and \(\det^{-k/2}\).

Since \(|r-t|<k\), subtracting the vacuum weights from
\((k/2,\ldots,k/2)\) gives positive entries on both compact
blocks of \(U(a-p,p)\). In the notation of Paul--Trapa's Lemma~3.2
the two joint-harmonic bounds are \(a-p\le r,\ p\le t\).
For the primed determinant all the entries after subtraction
are negative. The bounds for \(U(q,a'-q)\) become
\(a'-q\le r,\ q\le t\). They hold since \(r=(a-p)+(a'-q)\)
and \(t=p+q\). Thus each character occurs in the joint harmonics.
Part (1) of their Theorem~4.1 gives on the unprimed side
\(A_{\mathfrak q}(\det^{k/2}\boxtimes 1)\) with first block
\(U(a-p,p)\). Part (2) gives
\(A_{\mathfrak q}(1\boxtimes\det^{-k/2})\) with second block
\(U(a'-q,q)\) on the primed side. These are the same ordered
two-block parabolic before assigning its two inducing characters.
Finally the target splitting adds \(B/2\) to both block weights
on the first side and \(A/2\) on the second. Tensoring by this
determinant character commutes with cohomological induction, by
pulling the character through the induction and derived compact
invariant functors. Both modules therefore have the displayed
description. The shifts inherent in the authors' normalization
of \(A_{\mathfrak q}\) are identical on the two sides.
\cite[\S\S2.1, 2.3, 2.8, Lemma~3.2, Theorem~4.1 and Corollary~4.3]{paulTrapaHowe}

These are small theta quotients of the polynomial Fock modules,
hence irreducible nonzero targets of equivariant surjections.
Taking their algebraic tensor products over the real places gives
the quotient for the product Lie algebra and maximal compact.
Indeed the source polynomial module is spanned by finite products
of creation monomials from the separate places. A submodule of the
external product of simple Harish-Chandra modules is either zero
or the whole product, by the density theorem applied to finite
tuples in each factor (their commuting endomorphisms are scalars).
The quotient maps kill source derivatives and compact differences;
target derivatives and compact actions preserve matching. Extending
the maps linearly from polynomial vectors to their ambient
Schwartz spaces satisfies the stated real-map conditions, all of
which test polynomial inputs. For simultaneous archimedean
one-parameter groups one differentiates the finite product of
the individual actions, giving the sum of the placewise derivatives.
\end{proof}

\begin{theorem}[A nonzero supported undoubled period]\label{threal:nonzero}
For \(q=0\) and every allowed \(p\) the primed period is nonzero for some polynomial
Schwartz test satisfying the imposed full-rank condition at \(v_1\).
Both parities of \(a'\) are allowed.
\end{theorem}
\begin{proof}
Put \(Y=V'\), \(m=a'\), \(H=U(Y)\), \(J=U(W)\). The source quotient
is compact by the rank-zero case of reduction theory. Take probability
quotient measures for now; other Haar normalizations multiply the
periods by positive constants. We first prove the identity
\begin{equation}\label{threal:norm}
 \int_{[J]}\left|\int_{[H]}\Theta_\phi(g,h)\,dg\right|^2 dh
  = c_H\int_{H(\mathbb A_F)}
                \langle\omega(g,1)\phi,\phi\rangle\,dg,
                 \qquad 0<c_H<\infty .
\end{equation}
Here the scalar simply accommodates the choice of group Haar measure
on the right; it is one for quotient volume one with rational counting
measure. We use the usual self-dual Hilbert pairing, linear in the
first variable. For the proof it suffices to take finite pure tensors
with polynomial Fock factors at infinity.

On the doubled space \(Y\oplus(-Y)\), change from the product
polarization for \(\phi\otimes\bar\phi\) to the diagonal polarization.
The value at zero in the latter coordinates is
\(\int\phi(x)\overline{\phi(x)}\,dx\): the change is a difference and
average of the two position coordinates followed by Fourier transform
in the average variable. The Jacobian scalars can be fixed using
rational coordinates and globally self-dual measures, by the product
formula. At \((g_1,g_2)\) the analogous functional gives
\(\langle\omega(g_1,1)\phi,\omega(g_2,1)\phi\rangle\) times the
second-block determinant character of the doubled splitting. In the
product with the conjugate kernel that character is divided out.
Its exact conventional form is \(\chi_W(\det g_2)\) for the splitting
character on the \(H\)-action in the notation of the doubling
construction. Its modulus is one and it is trivial on rational points.
This block identity is also the change-of-model identity in
Gan--Qiu--Takeda, \S11.3, equations preceding (11.2)~\cite{ganQiuTakeda}.
The rational summation on the two product factors changes to the
doubled rational summation by Poisson summation.

For the given presentations, transport the standard implementations
using their rational Darboux coordinates. On each doubled factor the
possible discrepancy with the prescribed implementation on elements
supported at a single place is a scalar character trivial on the
Siegel stabilizer, by the argument of
Proposition~\ref{threal:inert}. This argument applies over the
real fields as well, using hyperbolic planes. At split places the
doubled group is a general linear group, generated by the block
parabolic and its conjugates by elementary elimination. Hence every
single-place discrepancy is trivial on each single factor by the
block rule. The global discrepancy is a continuous character:
compare the two implementations with a nonzero matrix coefficient
in a neighborhood of the identity, using unitarity and weak
continuity. It is trivial on the dense subgroup of elements supported
on finitely many places, hence everywhere. As a result the standard
automorphic theta and these periods agree after the rational change
of coordinates, up to the single scalar inherent in choosing the
unitary coordinate transform. Taking its modulus to be one suffices
in \eqref{threal:norm}.

To interchange the three integrations, use the chamber estimate of
Proposition~\ref{threal:decay} with \(W\) as integrating space of
dimension \(d\) and \(Y\) as partner of dimension \(m\). Its proof
applies to a skew-Hermitian integrating space as well. If a
Hermitian integrating space is used in the statement of the
Siegel--Weil formula, multiply its given skew-Hermitian form by
any \(0\ne\delta\in E\) with \(\bar\delta=-\delta\) and multiply
the other form by \(\delta^{-1}\). The product of the forms, the
underlying isometry groups, the flags and their moduli are preserved
in these conventions.
For \(g\) in compact source representatives the single kernel is
\(O(\prod T_i^{m/2})\), uniformly in \(g\). The squared bound is
integrable since the \(F\)-Witt index of \(W\) is at most \(p\)
and \(d-2p+1>m\). Now in the doubled Siegel--Weil formula integrate first over \([J]\),
with integrating space of dimension \(d\) and doubled group \(U(2m)\).
The parameters are \(s_0=(d-m)/2>m/2\), as \(a>a'\).
The coefficient of the probability theta integral is one by Weil's
convergent-range formula, applied with \(d>2m\); in the anisotropic
case this is also Theorem~7.1 of Gan--Qiu--Takeda. The Siegel Eisenstein
series with normalized induction is absolutely and locally uniformly
convergent in \(\Re s>m/2\); absolute convergence holds uniformly
on compact representatives of \([H]^2\) by majorization by the
positive spherical section at the real parameter. Thus integration
over \([H]^2\) permits absolute unfolding.
We use the probability normalization for the theta integral
\cite[introduction, Theorem~2 and \S3]{yamanaUnitary}.

Every maximal rational isotropic subspace of \(Y\oplus(-Y)\) projects
injectively onto each copy of \(Y\) since \(Y\) is anisotropic. Its
dimension is \(m\), so it is the graph of a rational isometry.
The set indexing the Siegel series is therefore one
\((H\times H)(F)\)-orbit with stabilizer \(\operatorname{diag} H(F)\).
After unfolding we integrate the matrix coefficient above, with its
second-block character cancelled, over
\(\operatorname{diag}H(F)\backslash H(\mathbb A_F)^2\).
The change \(g=g_2^{-1}g_1\) gives \eqref{threal:norm} and proves
absolute convergence of its right-hand integral. The same calculation
with any finite replacements of the test factors is valid.

We choose the factors so that the last integral is nonzero. At infinity
the primed source is negative definite. In the entry calculation of
Proposition~\ref{threal:compact} there are \(d-p\) negative planes
and \(p\) positive planes in each source row. Its vacuum determinant
power as a source representation is
\((A-B-d+2p)/2=j+p>0\). The \(d-p\) columns from the negative
planes are copies of the dual standard source representation.
Choose \(m\) of those columns, possible since \(d-p\ge m\), and
take their determinant to the power \(j+p\) times the vacuum.
This is a nonzero source-invariant polynomial vector. Its local matrix coefficient on
the compact source is the positive constant \(\|\phi_v\|^2\).
At a designated finite split place the conjugate polarization has
positions \(X\in M_{m,d}(F_v)\) and the source acts by left row
operations times a scalar character and the square root of the
modulus. Take the indicator of a sufficiently small compact open
ball about \((I_m\mid0)\). If \(g\) moves this ball to intersect
itself, the invertible leading block implies
\(g\in 1+\pi^b M_m(\mathcal O_v)\) for a prescribed \(b\);
this follows by writing \(g=X_1X_2^{-1}\) on the first blocks, or
its inverse if the action is inverse. Choose \(b\) so the character
is trivial there. The coefficient is nonnegative, vanishes off a
compact set and is positive on a neighborhood of \(1\).
This test satisfies the full-rank condition at \(v_1\).

At any specified inert place \(W\) has an integral hyperbolic
polarization by the lattice argument above, irrespective of the
parity of \(m\). The source acts on \(Y^{d/2}\) by linear changes
on all vectors times a scalar character. Now \(d/2>m\). Pick a small
ball about a tuple whose first \(m\) vectors are a basis of \(Y\).
The same leading-block argument inside \(\mathrm{GL}_m(E_{\mathfrak w})\)
forces overlap only in any prescribed neighborhood of \(1\) in
\(H(F_v)\). The scalar on this compact subgroup is trivial after
shrinking: a continuous character of a profinite group has open
kernel. The corresponding local integral again is strictly positive.

It remains to ensure that the unrestricted product of local
integrals is nonzero. Start with the factors just chosen at infinity
and \(v_1\), and the usual unit-norm lattice indicators at almost
all other places. There the hyperspecial group \(K_v\) fixes the
indicator: in the split linear model the unramified character is
trivial on units; in the inert hyperbolic polarization the source
acts by pullback with the unramified splitting. More explicitly the
character on the source for this polarization is trivial when the
partner is hyperbolic and its splitting character is trivial, as
is the case for the source splitting \(\alpha/\beta\) at inert places:
the two unramified characters have equal value on a uniformizer
by their common parity. The doubled Schr\"odinger model makes this
pullback implementation explicit, and the doubled uniqueness already
proved identifies it with the prescribed one. Take the local Haar
measures to give \(K_v\) mass one at almost all places. Write
\(b_v(g)=\langle\omega_v(g,1)\phi_v,\phi_v\rangle\) and
\(B_v=\int|b_v|\). Outside a finite set \(S\) one has \(b_v|_{K_v}=1\)
and \(B_v\ge1\). The integral of \(\prod_v|b_v|\) is finite by
the absolute unfolding, and is positive on some cylinder about
the identity. Exhaustion by finite restricted products and Tonelli
give \(\prod_{v\notin S} B_v<\infty\), hence
\(\sum_{v\notin S}(B_v-1)<\infty\). Thus
\[
        \left|\int_{H(F_v)} b_v(g)\,dg-1\right|
              \le B_v-1
\]
at these places. Only finitely many local integrals can vanish;
replace those and any exceptional finite factors by the compact
open tests just constructed. Absolute factorization now gives a
nonzero product. Formula~\eqref{threal:norm} yields a nonzero
undoubled primed period. The finitely modified tensor is a valid
supported test since the prescribed vacuum is standard almost
everywhere.
\end{proof}

\begin{theorem}[The scalar at the definite base]\label{threal:base}
Suppose \(p=q=0\). Both undoubled periods are nonzero, the primed
one on the required supported subspace. Common matching quotient
maps exist. For every choice of such maps there is a unique
\(c\in\mathbb C^\times\) such that \(I_\Phi=cI'_{\Phi'}\) for
all matched tests. Both parities are allowed. The inert and ordinary
split maps are furnished by Propositions~\ref{threal:allinert} and
\ref{threal:ordinarysplit}; the specified orbital maps at \(S\)
and the stronger kernel and support conditions at the totally split
good primes are those of Corollary~\ref{split:ram} and
Proposition~\ref{split:good} in the row polarizations. The
identification of those polarizations and their vacua is the one
in Proposition~\ref{threal:ordinarysplit}; at \(S\) use the
specified lattice bases. Proposition~\ref{threal:compact} gives
the archimedean maps.
\end{theorem}
\begin{proof}
First establish unprimed nonvanishing. Write \(n=a\),
\(H=U(V)\), \(J=U(W)\), \(\chi=\beta/\alpha=\nu^{-1}\),
\(T=S=S_0\sqcup Q\) the exact set above which \(\nu\) is ramified,
and \(s_0=(d-n)/2=a'/2\). Both adelic quotients are compact.
Form the Siegel section of \(U(V\oplus -V)\) through the
doubled Schwartz vector \(\phi\otimes\bar\phi\) at \(s_0\),
and extend it flat in normalized induction using the compact
subgroups specified below. Pair the resulting Eisenstein series
on the two copies of \([H]\) against the constants, cancelling the
second-block determinant character as above. Write this scalar
as \(Z(s)\). The graph orbit and the change of variables in
Theorem~\ref{threal:nonzero} express \(Z(s)\), for \(\Re s\)
large, as a positive measure constant times the product of the
local integrals of \(f_{v,s}(g,1)\) over \(H(F_v)\). Constants
are fixed by giving the hyperspecials volume one outside \(T\)
and choosing product Haar measure up to one global positive
factor. The meromorphically continued Eisenstein series can be
paired over the compact quotient: its meromorphic continuation
as smooth automorphic forms is locally uniform on compact sets
after multiplication by a local scalar denominator. At \(s_0\),
the anisotropic Siegel--Weil formula with integrating space \(W\)
and \(d>n\) gives \(Z(s_0)=C\|I_\phi\|^2\), \(C>0\), by Fubini.
No initial absolute convergence at \(s_0\) is used.
\cite[Theorem~7.1 and \S11.3]{ganQiuTakeda}

At each infinite place take a nonzero source-invariant polynomial
from Proposition~\ref{threal:compact}. Choose the maximal compact
in the doubled group containing \(H(F_v)^2\); hence
\(f_{v,s}(g,1)=\|\phi_v\|^2\) there. The doubled section is
compact-finite, since the doubled maximal compact at this place
is the product of the two compact groups. At \(v\in T\), which is
split because \(T\subset P\), take the small-ball test about a
full row frame in \(M_{n,d}(F_v)\) constructed above. On
\((g,1)\) the flat section equals the compactly supported matrix
coefficient at \(s_0\) times a positive Iwasawa height raised to
\(s-s_0\). Its integral is entire and strictly positive at \(s_0\).

At every other finite place the doubled section can be taken to
be the basic spherical section of value one at the identity.
At split places use the characteristic function of integral
matrices in the linear model. On the doubled row space the two
blocks together give the integral matrices with \(2n\) rows.
The determinant character is unramified and the integral general
linear group fixes their indicator. At inert places the extension
and characters are unramified, the additive character has conductor
\(\mathcal O_v\), and the residue characteristic is odd. The
lattices are perfect, \(W\) has even dimension and is integrally
hyperbolic, as proved above. In its polarization the doubled row
group acts linearly on \((V\oplus -V)^{d/2}\) with unramified
splitting, and the product of the two lattice indicators is
fixed by the stabilizer of the doubled self-dual lattice. This
works also for \(n\) odd. Indeed an integral hyperbolic plane
splits off repeatedly from a perfect Hermitian lattice of rank at
least two over an unramified quadratic extension of odd residue
characteristic by reduction and Hensel lifting, leaving a unit
line in odd rank. Thus the original row group too is unramified.
Normalize the indicators to have \(L^2\)-norm one. Evaluation on
the doubled diagonal then has value one. Transforming back to the
given presentation uses only finitely many departures from its
prescribed vacuum, exactly as above.

Let \(\eta=\eta_{E/F}\), and use a superscript \(T\) to omit finite
Euler factors over \(T\). The unramified doubling computation
with trivial representation of \(U(n)\) gives for the product of
the factors just described outside \(T\) and infinity
\[
 \frac{\displaystyle\prod_{i=0}^{n-1}
                 L^{T}_{E,\mathrm{fin}}(s+n/2-i,\chi)}
      {\displaystyle\prod_{i=0}^{n-1}
                 L^{T}_{F,\mathrm{fin}}(2s+n-i,\eta^{\,i})}.
\]
Here the base change of the spherical trivial representation has
the \(n\) exponents \((n-1)/2-i\): it is the spherical
constituent with trivial representation's Satake parameter,
the half sum of positive roots, restricted to the split group
over \(E\). In the inert odd case this list includes zero. The
doubling \(L\)-function uses \(s+1/2\), and \(\chi|_{\mathbb A_F^\times}=1\)
because the two ranks have equal parity. These give exactly the
shifts and the denominator above. See the unramified formula
\cite[\S4.2.1, (53)]{ehlsDoubling}. The product equals the
meromorphic continuation from its region of absolute convergence.
In particular no place with an unramified numerator has been
removed arbitrarily: at every prime of \(E\) over \(T\) the
character is ramified by the hypotheses, so its local factor is one.

At \(s_0\) the denominator arguments are at least \(a'+1>1\).
The numerator arguments are integers \(k=d/2-i\).
At each infinite place the angular exponent \(m\) of \(\chi\)
has \(|m|>d\); the gamma factors
\(\Gamma_{\mathbb C}(k+|m|/2)\) are finite and nonzero throughout
this interval. The Euler product gives nonvanishing when \(k>1\).
At \(1\), one elementary proof uses that both \(\chi\) and
\(\chi^2\) have nonzero infinity type. They have no Hecke pole.
For \(u>1\),
\(\zeta_E(u)^3|L_{E,\mathrm{fin}}(u,\chi)|^4
 |L_{E,\mathrm{fin}}(u,\chi^2)|\ge1\) off the ramified factors,
by \(3+4\cos t+\cos(2t)\ge0\) in the logarithmic Euler products.
The finite omitted factors of the three functions are bounded
above and below near \(1\). A zero would have order at least four
in this product, while \(\zeta_E^3\) has pole order three. For
\(k\le0\), apply the Hecke functional equation to the completed
function and use nonvanishing at \(1-k\) for \(\chi^{-1}\);
the conductor power and root number are nonzero
\cite{tateThesis}. This proves
\(Z(s_0)\ne0\). The primed supported period is already nonzero
by Theorem~\ref{threal:nonzero}.

We next compare the actual automorphic copies. On the definite
source the primed integral exists also for tests without the
support restriction. Use that extension in the following argument.
All three quotients are compact, with compact groups at infinity.
For any specified finite open compact level and infinite compact
type the space of automorphic vectors is finite-dimensional:
there are finitely many finite-adelic double cosets at that
level, and matrix coefficients of the prescribed compact type
span a finite-dimensional space on each. Consequently the unitary
spectrum is discrete with finite multiplicities. More explicitly
the orthogonal projection to any such level and type has finite
rank; these projections exhaust a dense subspace, and a continuous
direct integral of inequivalent unitary representations would have
infinite rank on any projection on which it is nonzero.

Let \(P_v,P'_v\) be the given local quotient maps and \(\pi\)
the restricted tensor product of the common irreducibles, using
their normalized vacuum images. At infinity take the entire
product of the compact factors; the polynomial calculation of
Proposition~\ref{threal:compact} gives its source-invariant
quotient with multiplicity one. For each of the two periods,
project to any irreducible automorphic summand of the target.
A nonzero projection, with all but one test factor fixed and
the other target factors paired with linear functionals, gives
an intertwiner from the local oscillator to the source trivial
representation times the local target constituent. At finite
places local Howe duality gives a unique such irreducible quotient
and a one-dimensional space of maps to it. This applies to
unitary dual pairs in characteristic zero and to the split
general linear dual pair; twisting an oscillator splitting by
characters preserves the assertion
\cite{ganTakedaHowe,minguezHowe}. At infinity the identical
conclusion follows from the compact calculation. Hence every
nonzero projection takes values in a copy of \(\pi\) and
annihilates each local kernel. Pairing is legitimate on algebraic
tensors: a smooth finite-level, compact-finite vector of an
irreducible automorphic summand belongs to its algebraic restricted
tensor product, and the pairings on the remaining factors separate
these vectors. At each finite tensor stage the kernel of the tensor
quotient is the sum of the subspaces having a kernel in one
factor, by exactness of tensor products over \(\mathbb C\).
This proves factorization of both periods. Smoothness then gives
pointwise equality whenever two resulting automorphic vectors are
equal in \(L^2\).

In fact \(\pi\) has multiplicity one here. Let \(i_0\) be the
unitary embedding defined by the nonzero primed image, using on
\(\pi\) the norm transported from that image. If another copy
exists, choose a unitary embedding \(i_1\) with image orthogonal
to \(i_0\). For any polynomial primed test \(u\) and compact-finite
smooth \(x\in\pi\), form the reverse lift
\[
  T_u(i_jx)(g)=\int_{[J]}(i_jx)(h)\,
                \overline{\Theta'_u(g,h)}\,dh
         \qquad (g\in U(V')(\mathbb A_F)).
\]
This function is constant in \(g\). In fact pairing against any
irreducible summand of the compact-quotient source spectrum and
varying \(x,u\) gives the same local intertwiners as above,
now with target constituent \(\pi_v\) fixed. Injectivity of
local Howe correspondence forces the source constituent to
be trivial at every place. At infinity the equivariant Hom in
the reverse direction also has this property: the Cauchy
decomposition in the compact calculation pairs distinct partitions
on the two compact groups. The globally trivial representation
occurs precisely on constant functions. Therefore \(T_u(i_1x)=0\):
its average on \([U(V')]\) is a positive normalization times
\(\langle i_1x,I'_u\rangle=0\).

On the other hand
\[
                 \|T_u(i_0x)\|_2^2=\|T_u(i_1x)\|_2^2.
\]
To see this form the flat doubled section \(f_{u,s}\) on
\(U(W\oplus -W)\) from \(u\otimes\bar u\), with parameter
\(s_*=(a'-d)/2=-a/2\). Pair its Eisenstein series, after the
second-block determinant correction, over \([J]^2\) against the
two copies of the vector \(i_jx\) and its conjugate. For \(\Re s\)
large there is just the graph orbit, since \(W\) is anisotropic.
Unfolding gives the integral over \(J(\mathbb A_F)\) of the
restricted section against the matrix coefficient of the right
regular action on \(i_jx\). In the convention with the section at
\((1,g)\) its scalar weight is \(f_{u,s}(1,g)\chi_*(\det g)^{-1}\),
where \(\chi_*=\alpha\) is the primed target splitting character.
The matrix coefficient is \(\langle R(g)i_jx,i_jx\rangle\), or
its conjugate upon conjugating the whole pairing. Either is
independent of \(j\) by isometry and equivariance. The two
meromorphic integrals thus agree identically. At \(s_*\) the
anisotropic Siegel--Weil formula is holomorphic, and its coefficient
is two for the probability theta integral because \(a'<d\).
The doubled product identity and Fubini over the compact quotients
give the two asserted squared norms times the same positive constant
\cite[Theorem~7.1 and \S11.3]{ganQiuTakeda}. Choose \(u\) with
\(I'_u\ne0\) and \(i_0x=I'_u\); then the average of
\(T_u(i_0x)\) is a positive multiple of \(\|I'_u\|^2\).
This contradicts \(T_u(i_1x)=0\). The two nonzero maps from the
common quotient now lie in the same one-dimensional intertwining
space by Schur's lemma. Restriction back to supported primed tests
and Corollary~\ref{threal:limits} give the required nonzero scalar
and its uniqueness.
\end{proof}

\begin{proposition}[The two periods against the cuspidal spectrum]\label{threal:orthA}
The unprimed period pairs to zero with every irreducible cuspidal
representation of \(J=U(W)\) whenever \(p>0\), for arbitrary \(q\).
If \(q=0<p\), the primed period also has this property, and on
removing a target hyperbolic plane the zero-vector orbit of the
unprimed line constant term vanishes identically.
\end{proposition}
\begin{proof}
Choose a rational prime completely split in \(E\) outside the finitely
many exceptional primes and the prime under \(v_1\), using the
splitting-prime theorem; let \(v\)
be a place above it. Write \(k=F_v\). By Lemma~\ref{split:ray}, the
unramified characters \(\alpha,\beta\) there have distinct values
on a uniformizer. The source-invariant quotient of the unprimed local
oscillator is the entire irreducible normalized induction
\[
       \pi_v=(\alpha\circ\det_a)\times(\beta\circ\det_{a'})
\]
by Proposition~\ref{split:good}. Thus every cuspidal constituent
pairing nontrivially with the unprimed period must have this local
factor. In detail the pairing, with other Schwartz coordinates and
other cuspidal coordinates fixed, gives an equivariant map into
the smooth dual of the conjugate of the local cusp representation.
Smoothness follows from invariance of the test under an open
compact subgroup, and the smooth dual in question is the original
representation by unitarity. Source invariance follows by changing
variables in the source integral.

For \(p>0\) the source has a rational hyperbolic plane: its form is
isotropic at every real place, it is split at the split finite places,
and at the inert places its perfect lattice has Witt index \(\lfloor
a/2\rfloor\). For the underlying \(F\)-quadratic norm form local
isotropy therefore gives global isotropy by Hasse--Minkowski, and
one completes the isotropic vector to a hyperbolic pair.
Write \(\Pi\) for such a cuspidal representation. Its adjoint lift
back to \(U(V)\), with conjugated oscillator, has a vector of
nonzero average, by Fubini. Indeed the absolute source theta sum in
Proposition~\ref{threal:decay}, integrated in \(g\), is bounded by
a power of \(H(h)\): in its unprimed estimate retain the factors
\(O(M^D H(h)^{A_B})\) before integrating the strictly negative
torus exponents. A cuspidal automorphic vector of finite compact
type decays faster than any power on a target Siegel set.
These estimates justify both the pairing and the interchange.

At any earlier source in the Witt tower, of dimension \(a-2b\)
with \(b\ge1\), no irreducible local representation can lift to
\(\pi_v\) in the forward oscillator. Indeed M\'inguez's split
correspondence adds the characters
\[
   \beta|\cdot|^{(1-l)/2},\
   \beta|\cdot|^{(3-l)/2},\ldots,\beta|\cdot|^{(l-1)/2},
               \qquad l=d-(a-2b)=a'+2b ,
\]
to the cuspidal support of the transformed source parameter.
The bare matrix oscillator adds the segment of the trivial
\(\mathrm{GL}_l\) representation. The scalar in the target row
action in \eqref{newtheta:big} is \(\beta(\det h)\) after removing
the unitary change-of-variables multiplier. The inverse transpose
on matrices interchanges the segment with its reversed dual, which
has the same multiset of real powers. This gives exactly the
displayed twist. The actual support of \(\pi_v\) has only \(a'\)
entries with unitary part \(\beta\). The other \(a\) entries have
unitary part \(\alpha\). This excludes the displayed segment.
\cite[Corollary~6.3(2)]{minguezHowe}
The first nonzero global adjoint lift of a cusp representation in a
Witt tower is cuspidal, by the tower constant-term formula
\cite[Proposition~3.1]{wuThetaTower}. If a lower lift were nonzero,
the first such lift would therefore have an irreducible cuspidal
constituent. Pairing against it and conjugating gives the excluded
forward local intertwiner. Thus the lift on \(U(V)\) is itself
cuspidal. The integral of a cusp representation against the constant
function on this positive-rank group is zero: a nonzero invariant
functional on a constituent would make it the trivial representation,
whose constant term along the given rational parabolic is nonzero.
This contradicts the preceding nonzero average.

For the primed period suppose \(q=0<p\) and let \(Y=V'\). A target
plane exists by the same norm-form argument, after multiplying its
skew form by a nonzero trace-zero element of \(E\).
The group \(U(Y)\) is anisotropic.
If the pairing with a cusp representation \(\Pi\) is nonzero, its
adjoint lift on \(U(Y)\) has nonzero projection to constants.
The entire adjoint lift is cuspidal by anisotropy, and is irreducible
by \cite[Theorem~5.3]{wuThetaTower}. Consequently it is the constant
representation. Its occurrence as functions consists of actual
constants, since right-invariant functions on \([U(Y)]\) are
constant. Wu's involution theorem now says that the forward lift of
these constants to \(W_1\) vanishes, where \(W_1\) is obtained from
\(W\) by deleting the rational hyperbolic plane
\cite[Theorem~5.1(3), Remark~5.2]{wuThetaTower}.
To check the twists in this use of the theorem, let \(d_X,d_Y\)
be its two splitting characters restricted to norm-one determinants
on the two groups \(X=W,Y\), respectively. Complex conjugation
of the oscillator replaces \(\psi,\chi_X,\chi_Y\) by their
inverses. Relative to Wu's oscillator for \(\psi^{-1}\) and the
original splitting characters it thus twists by
\(d_X^{-1}\boxtimes d_Y^{-1}\): on the doubled Siegel Levi this
follows from \(\chi(z/\bar z)=\chi(z)^2\); the doubled uniqueness
argument identifies the restrictions to the two factors. Wu's
reverse input is therefore \(d_X^{-1}\Pi\) and its output is
\(d_Y\) times the constants. The factor \(d_Y^{-1}\) in his
Theorem~5.1 cancels it in the forward lift.

This lower forward lift is nonzero by the calculation in
Theorem~\ref{threal:nonzero}. Here \(Y\) is unchanged, and the
new target has dimension \(d-2>2a'\), signature
\((d-p-1,p-1)\) with \(d-p-1\ge a'\), and is still hyperbolic
at the inert finite places, by Witt cancellation. The convergence
inequality there is \((d-2)-2(p-1)+1=d-2p+1>a'\).
These give respectively the absolute doubling identity, its real
source-invariant polynomial, and the required finite tests in that
proof. The splitting characters have unchanged parity on deleting
a plane. This contradiction proves the assertion for the primed
period as well.

Finally the zero-vector term on the unprimed side uses the
oscillator for the same \(V\) and the smaller \(W_1\). In the
polarization with position variable \(x\in V\) for the removed
target plane, setting \(x=0\) leaves the theta sum for
\(V\otimes W_1\); the extra source determinant character is
trivial for this added hyperbolic plane, whose source action on
the \(x\)-block is ordinary pullback. The source integral converges
by the proof of Proposition~\ref{threal:decay} with \(d-2\).
At a real place, the positive-character joint-harmonic condition in
Proposition~\ref{threal:allreal} now requires \(p\le p-1\).
Hence this local oscillator has zero trivial-source coinvariant
module: if its big theta module were nonzero, its unique irreducible
quotient would occur in the local correspondence. The source
integral at each \(h\) kills all the Lie algebra and compact-group
differences and so vanishes on the polynomial module. Continuity
and Hermite approximation give the same conclusion for all Schwartz
inputs. For the primed source the only rational isotropic vector is
zero, since \(Y\) is definite. Thus it has no nonzero-vector
orbit in this constant term.
\end{proof}

\begin{proposition}[Unfolding the larger nonzero orbit]\label{threal:coneA}
Let \(p>0\). Choose rational hyperbolic pairs \(x,y\) in \(V\)
and \(x',y'\) in \(W\), with orthogonal complements \(V_1,W_1\).
Write \(m(t,h_1)\) for the target line Levi and \(k=(a-a')/2>1\).
Use rational mixed Schr\"odinger coordinates and probability
measure on the target unipotent quotient. Let \(e_A\) be evaluation
in the coordinates dual to the modulation space
\[
        S_A= V\otimes_E Ex'\ \oplus\ Ex\otimes_E W_1.
\]
With \(K\) the source compact stabilizer and \(\rho\) the Heisenberg
action, put
\[
 T_A\Phi=\int_K\int_{\mathbb A_E^\times}
  e_A\rho(l\,x\otimes y')\omega(k_0,1)\Phi\
             \nu(l)|l|_E^k\,d^\times l\,dk_0 .
\]
The entire contribution of the rational nonzero isotropic vectors
to the unnormalized line constant term, at \(m(t,h_1)h\), is
\[
   c_A\,\alpha(t)|t|_E^{a'/2}\,
      I_{V_1,W_1}\big(T_A(\omega(1,h)\Phi)\big)(h_1).
\]
Here \(c_A\ne0\) depends on the rational models, the evaluations
and Haar measures; it is independent of the splitting characters.
At \(q=0<p\) this is the entire unprimed line constant term.
The entire primed constant term there is
\(\alpha(t)|t|_E^{a'/2}\) times its smaller zero-evaluation
period, with coefficient one for coherent mixed models and the same
source Haar measure.
\end{proposition}
\begin{proof}
First use \(V\otimes Ex'\) as modulation space. In its dual
coordinate the sum runs over \(z\in V(E)\). For fixed source and
target translates the sums are absolutely uniform on the compact
unipotent quotient. The center acts by an additive character of a
rational nonzero multiple of \(\langle z,z\rangle_V\); averaging
it leaves exactly the null vectors. The remaining target unipotent
coordinates have quotient \(W_1(\mathbb A_F)/W_1(F)\), writing
\(W_1\) here for its underlying \(F\)-vector group. They act on
the residual oscillator by the Heisenberg translations or
modulations \(z\otimes w\). This follows by applying the unipotent
matrix to the target coordinates \(x',W_1,y'\): the \(W_1\) shift
of \(y'\) is \(w\), and the quadratic correction in the \(x'\)
coordinate multiplies \(\langle z,z\rangle_V\), hence gives no
phase on this cone.

All nonzero rational null vectors form one \(U(V)(F)\)-orbit by
Witt extension. At the representative \(z=x\), use \(Ex\otimes W_1\)
as a further modulation space in the residual oscillator.
The last compact additive average kills every nonzero dual rational
coordinate by character orthogonality. Its surviving value is
therefore
\[
    \sum_{u\ {\rm rational}}
      (e_A\rho(x\otimes y')\,\omega(g,h)\Phi)(u),
\]
where \(u\) belongs to a position space for \(V_1\otimes W_1\).
The changes of rational polarization here preserve the theta
summation by Poisson summation with self-dual measures. Dual hyperbolic
bases may introduce a nonzero rational \(E\)-multiple of \(x\)
in this expression. Replacing the representative by that multiple
does not change the orbit, and its idele norm and Hecke character
are one.

Unfold the orbit to \(H_x(F)\backslash H(\mathbb A_F)\),
\(H=U(V)\), where the stabilizer of the vector is
\(H_x=N_x\rtimes U(V_1)\); the line stabilizer has the additional
Levi factor \(E^\times\). Write \(g=n\,m(b,g_1)k_0\) in source
Iwasawa coordinates. The adelic compact in this decomposition can
be taken to be the prescribed \(K\). At split finite places this
is integral row reduction. At an inert finite place one scales the
line generator to a primitive lattice vector and pairs it to one
by self-duality. The second vector can be made isotropic by the
surjectivity of the integral trace in the unramified extension.
Their unimodular hyperbolic plane splits off. Any two unitary
lattice complements here are isometric: a nonsingular finite-field
Hermitian space has a vector of nonzero norm, so integral
orthogonalization gives a unit diagonal, and every base-field unit
is an integral norm. The skew self-dual case is reduced to the
Hermitian case by the local different generator. These facts give
transitivity on isotropic lines. At infinity the two compact unitary
factors carry any isotropic line to the chosen line, by moving its
two components of equal length. The resulting local decompositions
are restricted products.

On the evaluated residual oscillator \(N_x\) acts trivially:
it fixes \(x\), acts with determinant one on \(S_A\), and acts
trivially on \(S_A^\perp/S_A=V_1\otimes W_1\). In mixed coordinates
the off-diagonal entries produce Heisenberg operators with parameters
linear in the dual \(S_A\)-coordinate, and quadratic phases in that
coordinate; both disappear at zero. It also fixes the displacement
\(x\otimes y'\). The source line Levi contributes
\(\nu(b)^{-1}|b|_E^{(d-2)/2}\) and the residual action of \(g_1\):
split off the whole target plane first, on which the source acts
by pullback on \(V\); on the remaining oscillator the source
Siegel character is \(\beta/\alpha\) and its modulation block
\(Ex\otimes W_1\) has dimension \(d-2\). Commuting the Levi left
past the displacement replaces \(x\otimes y'\) by
\(b^{-1}x\otimes y'\). The source Haar density on the line Levi
contains \(|b|_E^{1-a}\), because \(N_x\) has \(a-2\) ordinary
\(E\)-coordinates and one central \(F\)-coordinate of weight
\(b\bar b\). Integration over \([N_x]\) contributes its volume.
After the substitution \(l=b^{-1}\) the total character and modulus
are \(\nu(l)|l|_E^k\), since \(a-1-(d-2)/2=k\).
Integrating over \([U(V_1)]\) gives the smaller period.
The quotient integration constant and the unipotent volume are
positive for positive compatible Haar measures.

For justification of this unfolding one can integrate the
\emph{absolute value of the single orbit summand} over
\(H_x(F)\backslash H(\mathbb A_F)\) first. In the coordinates just
used its integral is bounded by the \(K,\mathbb A_E^\times\)
integral with weight \(|l|_E^k\) of the absolute smaller theta
sum integrated over \([U(V_1)]\). Choose the rational
source-adapted polarization of Proposition~\ref{threal:decay} for
this last sum. Its absolute integral is bounded by finitely many
Schwartz seminorms for each fixed finite test. The gap in dimension
is now \(k-1\) and the rational source Witt index is at most \(p-1\),
so the same negative-exponent estimates apply. In mixed coordinates
the new test is obtained by restriction to the linear coordinates
\((l,u)\); the partial Fourier transforms for changing the smaller
polarization act on \(u\) alone. At infinity these seminorms are
bounded by arbitrarily rapidly decreasing functions of \(l\),
uniformly on the compact \(K\)-orbit. At finitely many modified
finite places the restrictions are finite sums of products of
compactly supported tests in \(l,u\). At every other place they
are the indicator of the integral \(l\) times the smaller lattice
test. The absolute product is dominated by the Dedekind zeta
Euler product at \(k>1\); the finitely many local integrals at zero
also converge since \(k>0\). Tonelli applied to the translates of
the absolute orbit summand now justifies its global unfolding.
The same bounds prove that \(T_A\Phi\) is Schwartz and justify
interchanging its evaluation with the smaller period.
With multiplicative unit volume one at the basic places, the actual
scalar there is the absolutely convergent product of
\((1-\nu_w(\pi_w)(Nw)^{-k})^{-1}\), by the integral of the
lattice indicator along the displacement line.

On the target line Levi the evaluation scalar is
\(\beta(t)|t|_E^{a/2}\) because the whole \(V\otimes Ex'\)
block was removed first. Moving the Levi to the left past the
displacement changes \(l\) to \(l\bar t\); hence substitution in
the Mellin integral multiplies this scalar by
\(\nu(\bar t)^{-1}|t|_E^{-k}=\nu(t)|t|_E^{-k}\).
The target character is thus \(\alpha(t)|t|_E^{a'/2}\).
An evaluation with the same specified Heisenberg laws changes the
formula by a fixed scalar. Locally, invariance under the modulations
factors evaluation through restriction to zero: at a finite place
average on balls; at infinity differentiate the modulation law
to annihilate the coordinate multiples, which generate Schwartz
functions vanishing on that subspace (use Taylor's formula near it
and a cutoff away from it). The remaining intertwiner commutes with
all smaller translations and modulations, so is scalar. The
prescribed vacuum images fix the local normalizations at almost
every place. Thus changing the evaluations or the product measures
leaves the asserted independence from the splitting characters.
Finally Proposition~\ref{threal:orthA} removes the other unprimed
orbit when \(q=0\). On the primed side definiteness leaves only
zero, and the same first mixed evaluation gives the stated formula
with exponent \(a'/2\) and character \(\alpha\).
\end{proof}

\begin{proposition}[Detection by the entire line constant term]\label{threal:detectline}
Let \(f\) be a linear combination of the two periods of polynomial
tests, with the required support on the primed tests. In all
signatures, vanishing of the constant term of \(f\) along one
rational isotropic line, as a function on \(J(\mathbb A_F)\),
implies that \(f\) is cuspidal. For \(p+q>0\) this implies \(f=0\).
\end{proposition}
\begin{proof}
Witt extension conjugates any two rational isotropic subspaces of the
same dimension. If \(f\) is not cuspidal, choose the smallest \(r\)
for which its constant term along the stabilizer \(P_r\) of a rational
isotropic \(r\)-plane is nonzero. These are all the conjugacy classes
of maximal proper rational parabolics. By assumption \(r\ge2\).
On the Levi of \(P_r\), fix the other factor and vary the first
factor \(\mathrm{GL}_r(\mathbb A_E)\). Some resulting function,
allowing a compact right translate of \(f\) by Iwasawa decomposition,
is nonzero. It is cuspidal on this general linear group: the further
constant term along a \(j\)-plane in this block is the constant term
along the \((j,r)\)-flag, computed by unipotent Fubini, and this
also factors through the vanishing \(j\)-plane term for \(0<j<r\).

Here is the spectrum to which this cuspidal function belongs.
The original polynomial periods are compact-finite at infinity
and finite under the center of the target enveloping algebra:
source invariance factors the polynomial oscillator at each place
through its big theta module of the trivial representation, which
is of finite length by real Howe duality \cite{howeTranscending}.
Moderate growth follows from Proposition~\ref{threal:decay}.
The constant term restricted to its Levi is again center-finite.
Indeed on left unipotent-invariant functions restriction of a
central right differential operator gives its relative
Harish-Chandra projection (up to the modulus twist for the
unnormalized constant term). The Levi center is finite over this
image, since the invariants of the Levi Weyl group in the Cartan
polynomials are finite over the full Weyl invariants. This also
gives center-finiteness on the first direct factor of the Levi.

One can thus take a nonzero eigenfunction for the positive real
split center in the finite-dimensional span of its translates
(equivalently, extract an exponential coefficient on its Lie
algebra). These operations commute with finite-place translations.
Project also to a character on the compact norm-one idele class
center, and twist by a common real power of the idele determinant
to make the resulting central character unitary. Some such projection
is nonzero by compact Fourier decomposition. The resulting smooth,
compact-finite, center-finite cusp form is square-integrable modulo
the center by rapid decay on Siegel sets. It has a nonzero projection
to an irreducible unitary cuspidal representation of
\(\mathrm{GL}_r(\mathbb A_E)\). Its local factors are generic by
the cuspidal Whittaker expansion \cite{shalikaGeneric}. Thus at
each place above the \(v\) used in Proposition~\ref{threal:orthA}
the constant-term space must admit a generic irreducible subquotient
unitarizable after a determinant twist. Equivalently one can carry
out the center projections on the finite-place submodule generated
by the original slice: its split-center translates span a finite
number of copies of that module, and restriction to a central
eigenspace preserves its possible finite local subquotients.

At \(v\) each period factors through \(\pi_v\) by the full split
kernel computations, whether or not its coefficients at the other
places factor through their small theta quotients. The local
\(P_r\) has block sizes \((r,d-2r,r)\). Evaluation on the Levi after
taking the constant term factors through its unipotent coinvariants,
since a right unipotent translation becomes a translation of the
compact integration quotient. Normalizing the Jacquet functor only
twists each general linear block by a common real determinant power.
The geometric lemma computes the first \(\mathrm{GL}_r(k)\) action
in the resulting filtration, with possible subquotients of
\[
 \left(\alpha|\cdot|^{-(a-i)/2}\right)\circ\det_i
       \ \times\
 \left(\beta|\cdot|^{-(a'-j)/2}\right)\circ\det_j,\qquad i+j=r .
\]
Indeed a shuffle selects an initial block of length \(i\) of the
first determinant representation and of length \(j\) of the
second. The inverse square root of the modulus of their internal
parabolics gives the displayed exponents. The other end of the
isotropic plane accounts for the last, separate block of size \(r\).
\cite[geometric lemma, \S2.12]{bzgLemma}
The first block contributes a nondegenerate Whittaker functional
only if \(i,j\le1\): the highest derivative of the induced
representation is the tensor product of the highest derivatives
of its factors, and a determinant character of dimension \(>1\)
vanishes under that functor. The Whittaker coinvariant functor is
exact, so the same exclusion holds for every subquotient
\cite[\S4]{bzgLemma}.

It remains to examine \(r=2,\ i=j=1\). This normalized principal
series of \(\mathrm{GL}_2(k)\) is irreducible since its character
ratio has unitary part \(\alpha/\beta\) of infinite order and
therefore is neither \(|\cdot|\) nor \(|\cdot|^{-1}\).
The two Satake eigenvalues have distinct arguments and their
absolute values differ by the factor
\(|\pi_k|^{-(a-a')/2}\ne1\), where \(\pi_k\) is a uniformizer.
For unitarity the multiset must be invariant under
\(z\mapsto1/\bar z\), by the isomorphism to the conjugate dual.
That involution preserves the argument of each eigenvalue.
Consequently their distinct arguments force both absolute values
to be one. No common change of modulus can achieve this, giving the
required contradiction. Thus \(f\) is cuspidal. For \(p+q>0\)
Propositions~\ref{threal:orthA}, \ref{tpd:zero} and \ref{tpd:orth}
cover the two periods in every signature. Thus \(f\) is orthogonal
to the cuspidal spectrum and hence to itself; continuity gives
pointwise zero.
\end{proof}

\begin{corollary}[The detector in the \(a\)-descent]\label{threal:endA}
For \(q=0<p\), suppose that for one fixed \(c\) the
\(\alpha|\cdot|_E^{a'/2}\) components of the two periods agree up to
\(c\) on \(m(1,h_1)\) for every \(h_1\) and every polynomially
matched pair with supported primed test. Then those periods agree
up to \(c\) on all of \(J(\mathbb A_F)\) for every such pair.
\end{corollary}
\begin{proof}
By Proposition~\ref{threal:coneA} each entire line constant term
is just its indicated component. The matching condition survives
right translation by the target compact group: its real compact
action preserves the polynomial module, and its finite action
preserves full rank at \(v_1\) and intertwines the matching maps.
The hypothesis therefore applies also to \(m(1,h_1)k_0\) for each
target compact \(k_0\). The common Levi character extends this to
\(m(t,h_1)k_0\); left unipotent invariance and Iwasawa decomposition
cover the whole adelic group. Proposition~\ref{threal:detectline}
now applies to the difference.
\end{proof}

\begin{proposition}[The detector in the \(a'\)-descent]\label{threal:endB}
Let \(q>0\). Suppose a difference \(f\) of the two polynomial
periods, with supported primed test, has its full line constant
term transforming by \(\alpha(t)|t|_E^{a'/2}\) alone. Then \(f=0\).
In particular equality of the \(\beta(t)|t|_E^{a/2}\) components
of two matched periods, with one scalar for all the tests on
\(m(1,h_1)\), implies equality with that scalar on \(J(\mathbb A_F)\).
\end{proposition}
\begin{proof}
Induct on \(p\), keeping \(q>0\). For \(p=0\) the constant terms
of both periods have only the \(\beta\)-character by
Propositions~\ref{tpd:zero} and~\ref{tpd:unfold}. The two line
characters have distinct moduli since \(a>a'\), so the constant
term vanishes and Proposition~\ref{threal:detectline} applies.

For \(p>0\), split off planes in \(V,W\). At the real places one
can take compact groups adapted to the chosen planes. In detail,
choose on each hyperbolic plane a Cartan involution exchanging its
null lines and choose one on its orthogonal complement. The
resulting positive structure is conjugate to the original by a
real isometry. Translating the test by this isometry transports its
polynomial Fock module to the new one. Source translation leaves
the period fixed and simultaneous target translation replaces it
by a right translate, preserving purity and the zero assertion
to be proved. The lattices for the unused smaller finite models
can be chosen split along the rational decomposition at almost all
places, and compatible with the specified split polarizations at
\(v_1\). The induction uses only these rational theta models, their
finite tests, the prescribed signatures and the gap inequalities.
These hold in the smaller ranks \(a-2,a',d-2\), with indices
\(p-1,q\).

For each fixed target compact translate, restriction of this
pure constant term to the smaller unitary group is a difference
\(f_1\) of its two polynomial periods, with supported primed test.
Indeed the \(\alpha\)-component of the unprimed form is precisely
the cone in Proposition~\ref{threal:coneA}; the zero evaluation
has character \(\beta(t)|t|^{a/2}_E\). The smaller unprimed test
is polynomial at infinity by the orthogonal Fock factorization
of Lemma~\ref{tpd:mellin}, now with \(S_A\).
On the primed side the supported zero evaluation of
Lemma~\ref{tpd:so} gives precisely the \(\alpha\)-component.
The adjustment by source differences in that lemma may be made
anew for each fixed translate. Each resulting \(f_1\) in turn
has a two-character line constant term in its own ranks, by the
same calculations (and their definite-source versions). For an
arbitrary upper target translate, source Iwasawa is not needed
again: upper target Iwasawa and the pure character express its
slice as a scalar times a smaller target translate of a
compact-translate slice. Thus the same two line laws hold for all
slices used below.

We prove that the \(\beta\)-component of every such smaller
slice is zero. Choose any rational isotropic line in \(W_1\);
together with the first line it spans a rational isotropic
two-plane \(D\) in \(W\). For any fixed upper right translate
take the \(D\)-parabolic constant term of \(f\), and restrict its
first Levi block to \(\mathrm{SL}_2(\mathbb A_E)\), writing \(F\)
for the resulting automorphic function. Its ordinary upper
triangular constant term on this group is the average for the
isotropic \((1,2)\)-flag, by Fubini over the compact unipotent
quotients. Compute the same flag average in the other order by
first taking the outer line constant term and then the smaller
line term. These are averages along the two filtrations of the
\emph{same flag radical}. On \(d(t)=\operatorname{diag}(t,t^{-1})\)
the outcome, for each fixed right translate \(h\), is
\[
              F_N(d(t)h)=A(h)+\nu(t)|t|_E^{1-k} B(h),
                 \qquad k=(a-a')/2.
\]
The first outer factor is always \(\alpha(t)|t|_E^{a'/2}\).
The inner factors are \(\alpha(t^{-1})|t|_E^{-a'/2}\) and
\(\beta(t^{-1})|t|_E^{-(a-2)/2}\), respectively. These give
the two terms displayed. Their coefficients are continuous and
bounded on compact sets: they can be recovered by evaluating at
two fixed \(t\)'s with different idele norms.

In particular \(F\) belongs to
\(L^2(\mathrm{SL}_2(E)\backslash\mathrm{SL}_2(\mathbb A_E))\).
To check this, the displayed constant term is bounded on a Siegel
set since \(1-k<0\). Its difference from \(F\) decays faster than
any power of the split chamber parameter there. Indeed the
height exponent in Proposition~\ref{threal:decay} can be held fixed
for all right derivatives of the original period: choose once
the seminorm orders needed for integration; differentiating first
changes the test in these seminorms, not the powers of height in
the conjugation bounds. Integrating over the compact \(D\)-radical
therefore gives uniform moderate growth of derivatives of \(F\).
At fixed finite level the upper additive unipotent quotient is a
compact real torus. Subtract the mean of \(n\mapsto F(n d(t)c)\)
and use Fourier series there to bound its value by a Sobolev norm
of its positive-order derivatives of any sufficiently high order.
Along the archimedean positive chamber each unipotent derivative
gains the inverse positive root of \(d(t)\), as a right derivative
on \(F\). Compact factors \(c\) only change the fixed bounds. This
is the estimate in Corollary~\ref{tpd:detect} on the rank-one
group, with the bounded mean retained. Boundedness and finite
quotient volume give the assertion. This applies also to every
fixed finite local translate of the vector giving \(F\).

Use the good split place \(v\) of Proposition~\ref{threal:orthA}.
At one of its two \(E\)-places \(w\) the right
\(\mathrm{SL}_2(E_w)\) action on \(F\) factors through restriction
of the first block in the \(r=2\) Jacquet filtration computed in
Proposition~\ref{threal:detectline}. Its possible factors before
restriction are determinant characters and the one principal series
with \(i=j=1\). The latter has no unitarizable restriction
constituent. Here is a verification. Write \(G=\mathrm{GL}_2(E_w)\)
and \(G^+=Z(G)\mathrm{SL}_2(E_w)\). The quotient \(G/G^+\) is
the finite group of square classes. The irreducible principal
series \(\tau\) is finitely generated over \(G^+\) by coset
representatives applied to a generating vector. Choose a maximal
proper \(G^+\)-submodule; its finitely many \(G\)-conjugates
intersect trivially by \(G\)-irreducibility. The restriction
embeds in a finite sum of simple quotients, so is semisimple, and
its simple factors are \(G\)-conjugate. The center acts by a scalar
character. After twisting \(\tau\) by a real determinant power this
central character is unitary. If any factor on \(\mathrm{SL}_2(E_w)\)
had a positive invariant Hermitian form, so would every conjugate,
and their direct sum would give a \(G^+\)-invariant form.
Averaging it over \(G/G^+\) would unitarize the twist of \(\tau\),
contrary to the Satake calculation in Proposition~\ref{threal:detectline}.

The local subspace under consideration is a quotient of a
finite-length admissible \(\mathrm{SL}_2(E_w)\)-module: in the
Jacquet module take invariants of compact open subgroups fixing
the vector in all the other Levi blocks. Their invariant spaces
in the finite filtration are finite-dimensional on those blocks.
The first block restricts as just described. A finite-length
admissible submodule of a unitary Hilbert representation has
unitarizable simple constituents. In fact the closure of any of
its submodules has exactly that submodule's smooth vectors within
it, by compact-open averaging and finite-dimensionality at each
level; orthogonal projection then splits the submodule. Hence
only the trivial \(\mathrm{SL}_2(E_w)\) factors remain in the
space of translates of \(F\), and its local action is trivial.
In particular \(F_N(d(t))=F_N(1)\) if \(t\) is a uniformizer
at \(w\) and one elsewhere. Its modulus makes
\(\nu(t)|t|_E^{1-k}\ne1\). The flag formula now implies \(B(1)=0\).
The upper right translate was arbitrary, so the entire smaller
\(\beta\)-component vanishes.

The smaller slice \(f_1\) is thus pure for the \(\alpha\)-character.
Induction gives \(f_1=0\) for each upper compact translate.
The character law, upper unipotent invariance and Iwasawa
decomposition give vanishing of the full upper line constant term.
Proposition~\ref{threal:detectline} gives \(f=0\).
Finally equality on \(m(1,h_1)\) for all matched pairs kills
their entire \(\beta\)-component by compact stability of matching
and Iwasawa, as in Corollary~\ref{threal:endA}. The preceding
assertion applies to the remaining pure difference.
\end{proof}

\begin{proposition}[Matching all tests at an integral good split place]
\label{threal:localcone}
Let \(v\) lie above a completely split rational prime outside the
ramification set. Use unramified characters and a pair of integral
hyperbolic splittings in which the source, target and smaller basic
lattices decompose as direct sums and the displacement line is
primitive. Normalize the unit groups and \(K_v\) to volume one.
Write \(P_{1,A},P_{1,B}\) for the two smaller quotient maps,
normalized on their common vacuum. For every pair of locally
matched tests \(\phi,\phi'\), in the \(a\)-step one has
\[
   P_{1,A}(T_{A,v}\phi)=
    \prod_{w\mid v}L_w(\nu,k)\, P_{1,B}(e'_{0,v}\phi').
\]
In the \(a'\)-step the analogous equality is
\[
   P_{1,B}(T_{B,v}\phi')=
    \prod_{w\mid v}L_w(\nu,k+1)\, P_{1,A}(e_{0,v}\phi).
\]
Here the evaluations and basic vectors in the mixed coordinates
are normalized compatibly: evaluation of a basic vector is the
smaller basic vector, and its displacement evaluates to zero
outside the integral line. \(T_{B,v}\) is the local operator in
Lemma~\ref{tpd:mellin}. Both products range over the two places
of \(E\) above \(v\).
At a designated ramified split place the same pairs of comparison
maps span a space of dimension at most one in each direction.
\end{proposition}
\begin{proof}
Set \(K_0=F_v\). By Propositions~\ref{split:kernel} and
\ref{split:good}, or by Corollary~\ref{split:ram} at a designated
ramified place, the two source coinvariants are the same
irreducible \(\mathrm{GL}_d(K_0)\) representation \(\pi_v\).
In particular these are the precise kernels of any of the
surjective local maps in the matching data: a map killing the
source differences induces a nonzero map out of this irreducible
module. We show that the two maps to the smaller output after
applying a cone or zero evaluation factor through these
coinvariants.

For the zero evaluation, equivariance for the unchanged source
and its action on the smaller quotient give this immediately.
For \(A\), project
\(e_A\rho(x\otimes y')\omega(g,1)\phi\) to the smaller quotient.
It is left invariant by the full \(K_0\)-point vector stabilizer
\(H_x\), since its unipotent radical acts trivially as in
Proposition~\ref{threal:coneA}, and its remaining source factor
acts trivially after projection. This stabilizer is unimodular:
the action of the residual general linear block on the row and
column root groups has opposite determinant moduli, and is trivial
on the center. Integration over \(H_x\backslash H(K_0)\) therefore
gives a right \(H(K_0)\)-invariant map. Its Iwasawa coordinates
give precisely \(P_{1,A}T_{A,v}\), up to a nonzero Haar factor,
by the source Levi calculation above. The integral is absolute
since \(k>0\); the locally constant compactly supported restrictions
in the displacement and smaller position variables span a finite
number of smaller tests. For \(B\) use the vector stabilizer of
\(x\) in the primed source. It fixes the \(S_B\)-reduced quotient
and the class of \(y\otimes x'\) modulo \(S_B\), with no central
Heisenberg phase: the difference lies in
\((V'_1+Ex)\otimes Ex'\), isotropic together with
\(y\otimes x'\). The residual \(U(V'_1)\) again disappears after
projection. The Iwasawa calculation in Lemma~\ref{tpd:mellin}
gives \(P_{1,B}T_{B,v}\) by this quotient integral, now with
positive exponent \(k+1\).

Both resulting maps from \(\pi_v\) are invariant under the
unipotent radical of the target line parabolic. For \(A\), the
displacement changes there by an element of
\(Ex\otimes(W_1+Ex')\) in \(S_A\), with no Heisenberg phase
because \(x\) is isotropic. For \(B\), the displacement is fixed;
the target unipotent preserves \(S_B\), has determinant one there,
and acts trivially on the smaller tensor quotient. For the zero
evaluation the assertion holds already before the smaller quotient.
The actions on the reduced spaces in these checks are the mixed
ones at zero, as in Proposition~\ref{threal:coneA}. Moreover the
two maps in either comparison are equivariant for the smaller
target group and have the same line character, by the two Levi
computations.

At this split place the line parabolic has block sizes
\((1,d-2,1)\). Write its outer entries as \(u,z\); the element of
\(E_v^\times\) acting on the isotropic \(E_v\)-line is
\((u,z^{-1})\). The unnormalized line character in the \(a\)-step
is therefore \(\alpha(u)\alpha(z)|u/z|^{a'/2}\). Dividing by
\(\delta_P^{1/2}=|u/z|^{(d-1)/2}\) prescribes the first and last
characters of the normalized Jacquet module to be
\[
       \alpha|\cdot|^{-(a-1)/2},\qquad
       \alpha|\cdot|^{(a-1)/2}.
\]
The geometric lemma applied to the two determinant blocks of
\(\pi_v\) has exactly one such subquotient: the first and last
rows must both come from the \(\alpha\)-block, since
\(\alpha/\beta\) has infinite order on a uniformizer in the good
case, and is nontrivial on units in the ramified case. They leave
the normalized induction of the determinant characters of lengths
\(a-2,a'\) on the middle block. Its \(\alpha\)-exponent has no
shift because one row was removed at either end. It is irreducible
by the unitary induction theorem used in Proposition~\ref{split:good}.
In the \(a'\)-step the line character is \(\beta(t)|t|_E^{a/2}\);
the prescribed ends are \(\beta|\cdot|^{\mp(a'-1)/2}\). Both
come from the other block, leaving lengths \(a,a'-2\). Thus the
equivariant Hom into the stipulated smaller representation and
the two given characters has dimension at most one in each case,
by the Jacquet filtration and Schur's lemma. \cite[\S2.12]{bzgLemma}

At the stated good places evaluate the two maps on the upper
common spherical vector.
The zero evaluation gives the nonzero smaller vacuum. In the
cone integral the compact group fixes the test, and the shifted
evaluation gives exactly the smaller vacuum for integral displacement.
At each of the two places \(w\) the resulting geometric series
is \(L_w(\nu,k)\) or \(L_w(\nu,k+1)\), respectively. The product
is nonzero, fixing the scalar on every matched test.
\end{proof}

\begin{corollary}[The ramified unprimed comparison on all tests]\label{threal:ramA}
At a designated ramified split place, use the local character
\(\chi=\nu_w\), conductor \(\pi^c\), residue cardinality \(Q\)
and measures of Proposition~\ref{tpd:ram}. For all upper matched
tests the scalar in the \(a\)-step, multiplying the smaller primed
zero evaluation, is
\[
       Q^{c(k-1)}\chi(\pi)^c
       \sum_{r\in(\mathcal O/\pi^c)^\times}
                         \chi(r)^{-1}\psi_v(r/\pi^c).
\]
\end{corollary}
\begin{proof}
Use the pivot-determinant tests of Proposition~\ref{tpd:ram},
now removing the first and last of the \(a\) target pivot columns
on the unprimed side. The upper and lower pairs are matched by the
same integral row computation there, and both orbital functions
have value one at the identity before applying the reductions.
The primed zero evaluation integrates one nonpivot column with
\(a'\) entries in \(\pi^c\mathcal O\), giving
\(Q^{-a'c}\phi'_1\). On the unprimed side write
\(l=(s,t)\) in the two split components. The shift \(l x\otimes y'\)
has position entry \(s\) at the first row and last pivot column
and momentum entry \(-t\) at the last row and first pivot column:
in the chosen idempotent the coordinates of \(x,y'\) are the
first and last basis vectors; pairing their conjugates into the
dual basis gives the last and first vectors, with the minus sign
from the skew target pair. The two entries have zero mutual
Heisenberg half-pairing.
\src{src/Development/WeilBound/ThetaLifts/LocalTheory/StandardValuations/AdelicInclusionMapPlaceSection.lean}{218}{Paired-coordinate identities for the unprimed displacement}.

In the \(S_A\)-evaluation one integrates the first pivot column
and the first source row except for the last pivot column. This
follows by using the translations in \(V\otimes Ex'\) first
and then those in \(Ex\otimes W_1\); the dual modulation
coordinates are set equal to zero. Thus the pivot matrix has
the form
\[
      \begin{pmatrix} *&*&s\\ *&D_1&0\\ u&0&0\end{pmatrix}
\]
with Fourier factor \(\psi_v(-tu)\). Its determinant is
\(-su\det D_1\). Apart from \(u\), the integrated pivot entries
have unit volume, and the \(a'\) integrated nonpivot entries of
the first row give \(Q^{-a'c}\). Replacing \(u\) by \(-u\)
removes the two minus signs. With
\(D_\chi=\mathbf1_{\mathcal O^\times}\chi^{-1}\), this yields
\(Q^{-a'c}D_\chi(s)\widehat D_\chi(t)\phi_1\).
The compact source average fixes the test. The Mellin weight is
\(\chi(s)\chi(t)^{-1}|st|^k\). The \(s\)-integral is one; the
\(t\)-integral is supported on \(t=\pi^{-c}\mathcal O^\times\),
where \(\widehat D_\chi(\pi^{-c}u)=
 Q^{-c}\chi(u)G(\chi,\psi_v)\) by the unit residue sum.
Consequently this integral is \(Q^{c(k-1)}\chi(\pi)^c
G(\chi,\psi_v)\). Cancel the common factor \(Q^{-a'c}\) and use
Proposition~\ref{threal:localcone} to pass from this nonzero
smaller vector to all tests.
\end{proof}

For an adapted integral plane let \(\mathfrak b x\) and
\(\mathfrak c x'\) be its two first-line lattices. The duality
condition on the skew lattice, using the absolute trace, gives
the second target line lattice
\(\mathfrak D_{E/\mathbb Q}^{-1}\bar{\mathfrak c}^{-1}y'\).
Indeed a coefficient must pair with every element of
\(\mathfrak c x'\) into the codifferent, and the other integral
summands are orthogonal. Thus the finite displacement
\(l x\otimes y'\) is integral precisely for \(l\) in
\(\mathfrak b\mathfrak D_{E/\mathbb Q}^{-1}
  \bar{\mathfrak c}^{-1}\widehat{\mathcal O}_E\).
These ideals are units at the specified ramified places by the
prescribed paired integral bases and total rational splitting.
The relative different for \(E/F\) is trivial and
\(\nu|_{\mathbb A_F^\times}=1\), by the common parity of the
two splitting characters. Hence the character of the codifferent
is one, and that of the displacement ideal is
\(\nu(\mathfrak b)\nu(\mathfrak c)\); the norm factor with
exponent \(-k\) depends only on the chosen lattices. This gives
the two positive ideal exponents for the \(a\)-step. In particular
the splitting-prime calculation of the two comparison factors also
allows nonprimitive line generators by the corresponding change of
variable in the local Mellin test of the lattice vacuum.
\src{src/Development/WeilBound/ThetaLifts/LocalTheory/ConeShifts/IdealId0.lean}{49}{The two displacement ideals and their character values}.

\paragraph{The scalar comparison for \(p+q>0\).}
The common local maps give an algebraic test for the
global discrepancy. Fix polynomial real and finite tests and take the
restricted tensor products with the specified vacuum lines. Let
\(P,P'\) be the tensor maps to the common quotient, using the
full-rank restriction on the primed domain. Suppose the quotient
and full-rank surjectivities in \S\ref{newtheta:global} have been
established at the places in question. Then \(P,P'\) are onto:
each pure tensor has finitely many non-vacuum factors; choose lifts
factor by factor and use the prescribed vacua elsewhere. Work with
the rational summation functional averaged over the source, denoted
by \(I_h,I'_h\). For fixed \(h\) its action on source differences
vanishes by change of variables wherever the summation descends to
the source quotient. The first additional global requirement is
\[
            I_h(\ker P)=0,\qquad I'_h(\ker P')=0 .
\]
In particular a local quotient that kills the source differences
does not identify the annihilator of the global integral unless its
extra kernel is killed too. Once these equalities hold, the two
periods give well-defined linear maps \(A,A'\) from the common tensor
module to functions of \(h\). Proportionality is equivalent, when
\(A'\ne0\), to the following directly testable equalities for all
module vectors \(x,y\) and adelic \(h,h_0\):
\[
                   A(x)(h)\,A'(y)(h_0)
                    =A'(x)(h)\,A(y)(h_0).
\]
Choose \(y,h_0\) with \(A'(y)(h_0)\ne0\) to prove this equivalence;
the scalar is nonzero exactly when \(A\ne0\). Corollary~\ref{threal:limits}
then gives the assertion on the entire matching space. Conversely a
nonzero primed period on some supported Schwartz test already gives
one on a polynomial test: keep its finite factors fixed, approximate
each archimedean Schwartz coefficient by its finite Hermite
expansions, and use the continuity proved above. Hermite expansions
converge in all Schwartz seminorms, by iterating the harmonic
oscillator on the coefficients and using the creation and annihilation
operators to control the seminorms of its eigenfunctions.

Theorem~\ref{threal:base} establishes the scalar comparison at the
definite base. Theorem~\ref{threal:nonzero} proves \(A'\ne0\) also
for \(q=0<p\) whenever the factorization exists, and gives a nonzero
supported period there before factorization. The graph-orbit proof
of the base comparison changes once \(W\) has isotropic lines:
maximal isotropic subspaces of \(W\oplus -W\) can project with
nonzero kernels, and their terms survive in the unfolding unless
one proves the necessary vanishing. Proposition~\ref{threal:allreal}
gives the real quotient in these signatures, and
Proposition~\ref{threal:detectline} reduces the global comparison to
one entire line constant term for all target translates.

The rank restriction on other available multiplicity formulas
matters here. The determinant characters on the two inducing blocks
give the two Satake strings of lengths \(a,a'>1\) at the completely
split good primes. The natural Arthur parameter with these strings
has two character blocks with nontrivial Arthur \(\mathrm{SL}_2\);
the generic multiplicity-free formula does not apply to it.
Chen--Zou's Theorem~7.7 treats elliptic parameters at \(F\)-rank
one and also those for which each nontrivial Arthur block has
cuspidal degree greater than the rational Witt index. Here those
degrees are one. Moreover this theorem describes a discrete
spectrum in terms of packets, defined as multisets; using it on a
specific period would require proving discrete realization and
identifying its packet member. Even a multiplicity of one does not
give the ratio of functionals in different ranks.
\cite[\S5.3, Assumption~7.2 and Theorem~7.7]{chenZouMultiplicity}

In the descent in \(a\), the actual comparison of tests must send
the larger cone integral to the smaller zero evaluation times
\(C_A\,\mathrm{stepFac}_A\,L^{S_0}(\nu,(a-a')/2)\).
Proposition~\ref{threal:coneA} identifies the two global constant
terms with those operations and Corollary~\ref{threal:endA} proves
the required detection. Identifying their images in the smaller
matching quotients with the uniform scalar is a separate assertion.
In the descent in \(a'\), the smaller cone integral must match the
larger zero evaluation times
\(C_B\,\mathrm{stepFac}_B\,L^{S_0}(\nu,(a-a')/2+1)\), and must
land in the smaller supported space. Lemma~\ref{tpd:mellin}
proves that support assertion, Proposition~\ref{tpd:unfold}
identifies the cone component, and Proposition~\ref{threal:endB}
proves detection by it in every signature. The signs in the two
finite products and the conductors are as in \S\ref{newtheta:global}.
The constants \(C_A,C_B\) precede the quantifiers over the characters,
both Weil systems, and both sets of finite quotient maps once the
evaluations, real maps and measures are fixed.
Proposition~\ref{threal:localcone} determines both ratios on all
tests at the integral good split places in its hypotheses.
Proposition~\ref{tpd:ordinary} extends these identities to every
unramified place split in \(E/F\), including resonant characters
and the factors from the adapted line lattices.
Corollary~\ref{threal:ramA} and Proposition~\ref{tpd:ram} determine
the ratios at the designated ramified split places.
At inert places Proposition~\ref{threal:allinert} gives the
irreducible quotient and its common vacuum. Propositions~\ref{tcr:descend} and \ref{tcr:unique} prove that the
cone and zero operations factor through the upper quotient on every
local test and compute their ratio after projecting to the smaller quotient. At the real places
Proposition~\ref{threal:allreal} identifies the two modules in
each fixed rank. Proposition~\ref{tcr:unique} compares the compact-averaged
Mellin operators and evaluations between the two ranks on the
polynomial vectors for the fixed real quotient maps. Theorem~\ref{tcr:global}
assembles the comparisons for the prescribed polarizations, lattice vacua
and product measures, giving the uniform constants and nonvanishing in
positive signature.
\src{src/Development/WeilBound/ThetaLifts/ConeHypotheses/PolDense.lean}{46}{The two normalized cone comparisons};
\src{src/Development/WeilBound/ThetaLifts/ConeHypotheses/PolDense.lean}{32}{Detection by the selected constant-term component}.
